\documentclass[11pt,a4paper]{article}
\usepackage{fontspec}
\setmainfont{NotoSerifCJKkr-Regular.otf}[
  Path=assets/,
  BoldFont=NotoSerifCJKkr-Regular.otf,
  ItalicFont=NotoSerifCJKkr-Regular.otf,
  BoldItalicFont=NotoSerifCJKkr-Regular.otf]
\setmonofont{Latin Modern Mono}[Scale=MatchLowercase]
\XeTeXlinebreaklocale "ko"
\XeTeXlinebreakskip=0pt plus 0.08em
\usepackage{amsmath,amssymb,amsthm,mathtools}
\usepackage{geometry,graphicx,booktabs,array,longtable,tabularx}
\usepackage{setspace,titlesec,caption,enumitem,placeins,needspace}
\usepackage{tikz,pgfplots}
\usetikzlibrary{arrows.meta,positioning,calc}
\pgfplotsset{compat=1.18}
\usepackage[numbers,sort&compress]{natbib}
\usepackage{xurl}
\usepackage[unicode,hidelinks,pdfencoding=auto]{hyperref}
\geometry{left=26mm,right=26mm,top=25mm,bottom=25mm,footskip=12mm}
\setstretch{1.18}
\setlength{\parindent}{1em}
\setlength{\parskip}{0.15em}
\setlength{\emergencystretch}{1.5em}
\setlength{\tabcolsep}{4pt}
\renewcommand{\arraystretch}{1.16}
\setlength{\abovedisplayskip}{9pt plus 2pt minus 2pt}
\setlength{\belowdisplayskip}{9pt plus 2pt minus 2pt}
\setlength{\textfloatsep}{13pt plus 3pt minus 2pt}
\setlength{\intextsep}{13pt plus 3pt minus 2pt}
\setcounter{topnumber}{3}
\setcounter{bottomnumber}{2}
\renewcommand{\topfraction}{0.90}
\renewcommand{\bottomfraction}{0.80}
\renewcommand{\textfraction}{0.08}
\renewcommand{\floatpagefraction}{0.72}
\titleformat{\section}{\normalfont\large}{\thesection.}{0.55em}{}
\titleformat{\subsection}{\normalfont\normalsize}{\thesubsection.}{0.55em}{}
\titleformat{\subsubsection}{\normalfont\normalsize}{\thesubsubsection.}{0.55em}{}
\titlespacing*{\section}{0pt}{2.0ex plus .7ex}{1.0ex plus .3ex}
\titlespacing*{\subsection}{0pt}{1.6ex plus .5ex}{0.6ex plus .2ex}
\captionsetup{font=small,labelfont=normalfont,labelsep=period,skip=6pt}
\captionsetup[table]{position=top}
\renewcommand{\abstractname}{초록}
\renewcommand{\figurename}{그림}
\renewcommand{\tablename}{표}
\renewcommand{\refname}{참고문헌}
\renewcommand{\proofname}{증명}
\setlist{nosep,leftmargin=1.8em,topsep=3pt}
\newcolumntype{P}[1]{>{\raggedright\arraybackslash}p{#1}}
\newcolumntype{Y}{>{\raggedright\arraybackslash}X}
\newtheoremstyle{paper}{6pt}{6pt}{\normalfont}{0pt}{\normalfont}{.}{0.5em}{}
\theoremstyle{paper}
\newtheorem{definition}{정의}[section]
\newtheorem{proposition}{명제}[section]
\newtheorem{contract}{계약}[section]
\numberwithin{equation}{section}
\newcommand{\Stream}{\operatorname{Stream}}
\newcommand{\Val}{\operatorname{Val}}
\newcommand{\cut}{\operatorname{cut}}
\newcommand{\Dep}{\operatorname{Dep}}
\newcommand{\Anc}{\operatorname{Anc}}
\newcommand{\tr}{\operatorname{tr}}
\newcommand{\diag}{\operatorname{diag}}
\newcommand{\rank}{\operatorname{rank}}
\newcommand{\Live}{\mathsf{Live}}
\newcommand{\Record}{\mathsf{Record}}
\newcommand{\code}[1]{\texttt{#1}}
\newcommand{\den}[1]{\mathopen{[\![}#1\mathclose{]\!]}}
\hypersetup{pdftitle={다중 하드웨어 BCI를 위한 Rust 기반 도메인 특화 언어의 형식 설계},
  pdfauthor={},pdfsubject={언어 의미론, 신호 처리, 하드웨어 계약, 시간 동기화, Rust 실행 구조},
  pdfkeywords={BCI, DSL, Rust, formal semantics, hardware, clock, provenance}}
\begin{document}
\begin{center}
{\fontsize{16}{23}\selectfont 다중 하드웨어 BCI를 위한 Rust 기반\\도메인 특화 언어의 형식 설계\par}
\vspace{7pt}
{\small 설계 명세 1.0\qquad 2026년 10월 31일}
\end{center}
\begin{abstract}
본 논문은 전기생리와 광학 뇌신호를 처리하는 Rust 기반 도메인 특화 언어의 수학적 구조를 제시한다.
언어는 측정량, 단위, 채널 스키마, 시계 영역, 표본 발생 규칙을 스트림 형식에 포함하고,
측정 시각과 가용 시각을 함께 사용하여 실시간 연산의 인과성을 정의한다.
상태 전이 연산자의 직렬ㆍ병렬 합성, 초기 지연을 갖는 피드백, 학습ㆍ평가 데이터 계보를
하나의 형식 판단으로 연결한다. 하드웨어 계층은 장치 모델, 펌웨어, 획득 경로와 실행 환경을
지원 프로필로 묶고, 구성별 기능 집합과 유효 설정 공간으로 연결 조건을 계산한다.
다중 장치의 시계 변환에는 유효 구간과 오차 반경을 부여하며, 창 연산은 입력 진행 경계로 완결 시점을 결정한다.
실행 구조는 소유권이 명시된 장치 actor, 불변 데이터 블록, 모델 활성화 사건과 유한 버퍼 실행 계획으로 구성한다.
논문은 합성의 형식ㆍ인과성 보존과 지연 피드백의 실행 유일성을 증명하고,
EEG 및 고밀도 전기생리의 수치 예제, 18개 하드웨어 프로필, 규격 적합성 조건을 제시한다.
\end{abstract}
\noindent{\small 주제어: 뇌--컴퓨터 인터페이스, 도메인 특화 언어, 형식 의미론, 시계 동기화, 데이터 계보, Rust}
\section{설계 목적과 연구 기반}
\label{sec:introduction}
뇌--컴퓨터 인터페이스(BCI)는 신호 획득, 신호 변환, 상태 추정, 사건 생성과 피드백을 연결한다.
하나의 실험은 전극 전압과 광학 측정, 입력 사건과 모델 상태, 여러 장치의 시계를 함께 다룬다.
본 설계는 이 관계들을 언어의 형식과 실행 의미에 포함한다. 프로그램은 필요한 측정량과
시간 관계를 기술하고, 컴파일러는 선택한 하드웨어 구성과 연산자의 계약을 결합하여 실행 계획을 생성한다.

BCI2000은 신호 획득, 신호 처리, 응용 기능을 모듈로 구성하고 실험 조건과 상태를 공유하는
구조를 제시하였다\cite{bci2000}. BCILAB은 신호 처리와 학습 절차를 조합하고 방법을 평가하는
프레임워크를 구성하였다\cite{bcilab}. MOABB는 데이터셋, 전처리와 평가 인터페이스를 통합하여
BCI 알고리즘 비교의 재현성을 다룬다\cite{moabb}. 본 논문은 이 세 작업에서 사용하는
획득 경계, 처리 그래프, 평가 단위를 명시적 언어 요소로 연결한다.

형식적 기반은 스트림 의미론, 동기 데이터흐름, 소유권에 둔다. Kahn의 프로세스 네트워크는
통신 이력을 접두 순서와 함수 합성으로 해석한다\cite{kahn}. Synchronous Data Flow(SDF)는
연산자의 생산ㆍ소비 토큰 수를 이용하여 다중 표본률 실행 계획을 구성한다\cite{sdf}.
Lustre는 논리 시계와 지연을 반응형 프로그램의 구조에 포함한다\cite{lustre}.
RustBelt는 소유권과 차용의 의미 및 라이브러리 경계의 검증 조건을 제시한다\cite{rustbelt}.
본 설계는 스트림에 측정 시각과 가용 시각을 부착하고, 각 가변 자원에 소유 actor를 배정하며,
시계ㆍ학습ㆍ하드웨어 계약을 그래프의 합성 조건으로 사용한다.

설계의 주요 산출물은 다섯 가지이다. 첫째, 물리 차원과 명목 채널 의미를 포함한 스트림 형식이다.
둘째, 상태 전이와 가용 시각 인과성을 결합한 코어 의미론이다. 셋째, 구성 공간과 기능 집합을
사용하는 하드웨어 바인딩이다. 넷째, 시계 변환ㆍ창 완결ㆍ버퍼 용량의 계산 규칙이다.
다섯째, 모델 계보와 실행 사건을 재생 가능한 기록으로 묶는 Rust 실행 구조이다.

\subsection{정의, 근거, 계산의 구성}
논문은 정의와 계약으로 언어의 동작을 지정하고, 인용 문헌으로 장치 규격과 선행 연구의
근거를 연결한다. 명제는 지정한 코어 의미론에서 합성 성질을 증명한다. 수치 예제는
장치 규격과 명시한 실행 매개변수를 식에 대입하여 처리량ㆍ지연ㆍ용량을 계산한다.
각 하드웨어 프로필은 원문의 기관, 제목, 버전 또는 열람일을 함께 가진다.
2026년 10월 31일의 문서 판본을 참고문헌에 기록한다.

\begin{figure}[tb]
\centering
\begin{tikzpicture}[>=Latex,font=\small,node distance=15mm]
\node (source) at (0,1.4) {프로그램과 프로필};
\node (ir) at (4.6,1.4) {형식화 중간 표현};
\node (exec) at (9.2,1.4) {Rust 실행 계획};
\node (contract) at (0,-0.3) {측정ㆍ시간ㆍ기능 계약};
\node (sem) at (4.6,-0.3) {스트림 의미와 불변식};
\node (trace) at (9.2,-0.3) {값ㆍ사건ㆍ계보 기록};
\draw[->] (source)--node[above]{구체화}(ir);
\draw[->] (ir)--node[above]{계획 생성}(exec);
\draw[->] (source)--(contract);
\draw[->] (ir)--(sem);
\draw[->] (exec)--(trace);
\draw[->] (contract)--node[below]{합성 조건}(sem);
\draw[->] (sem)--node[below]{실행 관계}(trace);
\end{tikzpicture}
\caption{언어 구체화와 실행 기록의 대응. 위쪽 경로는 컴파일 과정을, 아래쪽 경로는 각 산출물에 붙는 의미와 계약을 나타낸다.}
\label{fig:architecture}
\end{figure}

\subsection{언어의 조직}
표면 언어는 장치 프로필 선택, 채널 스키마, 신호 연산, 학습 프로토콜, 시간 계약을 선언한다.
컴파일러는 유리수 표본률, 정수 창 길이, 유한 채널 집합과 명목 식별자를 정규화한다.
구성 조회로 얻는 값은 런타임 기술자와 연결한다. 이후 형식화 중간 표현은 원시 연산자,
직렬ㆍ병렬 합성, 지연, 피드백으로 구성한다. 이 조직은 새로운 장치와 알고리즘을
프로필과 원시 연산자 계약으로 추가하는 확장 지점을 제공한다.

측정값의 물리 단위, 장치 시각, 모델의 데이터 계보는 각각 독립적인 정보이다.
언어는 이 정보를 동일한 처리 그래프 안에 배치하고 변환마다 입력ㆍ출력 관계를 기록한다.
예를 들어 전압 단위 변환은 수치 배율을, 재기준화는 채널 스키마를,
시계 정렬은 시간 뷰를, 모델 갱신은 매개변수 버전과 가용 시각을 변경한다.

\section{측정 데이터와 스트림 형식}
\label{sec:data}
\subsection{물리 차원, 단위와 명목 측정량}
물리 차원을 SI 기본량의 지수 벡터로 표현한다.
\begin{equation}
 d=(d_M,d_L,d_T,d_I,d_\Theta,d_N,d_J)\in\mathbb Z^7.
\end{equation}
예를 들어 전압의 차원은 $(1,2,-3,-1,0,0,0)$, 주파수의 차원은
$(0,0,-1,0,0,0,0)$이다. 곱셈은 차원 벡터를 더하고 나눗셈은 차원 벡터를 뺀다.
단위 레코드 $u=(d,s,o)$는 기준 단위와의 관계 $x_{\mathrm{SI}}=sx+o$를 제공한다.
전기생리 전압과 농도 변화는 각 연산이 사용하는 선형 단위 표현을 선언한다.

채널 $c$의 기술자를 다음과 같이 정의한다.
\begin{equation}
 \chi_c=(\mathrm{id}_c,q_c,d_c,u_c,r_c,g_c,\mathrm{role}_c),
 \qquad \mathcal C=(\chi_1,\ldots,\chi_C).
 \label{eq:channels}
\end{equation}
$q_c$는 명목 측정량, $r_c$는 기준, $g_c$는 위치ㆍ기하 정보이다.
전극 전압의 식별자는 측정 전극과 기준 전극 또는 기준 연산을 포함한다.
광학 측정의 식별자는 광원, 검출기, 파장과 측정 종류를 포함한다.
기하 정보는 좌표값과 함께 좌표계 식별자, 길이 단위, 장치 배치를 보관한다.

\begin{table}[tb]
\centering\small
\caption{채널 역할별 의미와 형식 정보}
\begin{tabularx}{\linewidth}{@{}P{0.19\linewidth}YY@{}}
\toprule
역할 & 측정량과 단위 & 식별 정보 \\
\midrule
전기생리 & 전극 전압, V 또는 $\mu$V & 전극, 기준, 위치, 증폭 설정 \\
광학 원시값 & 광강도에 대응하는 장치 단위 & 광원, 검출기, 파장, 스케일 \\
광학 파생값 & 광밀도 또는 농도 변화 & 측정 종류, 단위, 광경로 모형 \\
운동 센서 & 가속도, 각속도 & 축, 좌표계, 센서 위치 \\
디지털 사건 & 상태 변화, trigger word & 입출력 선, 극성, 관측 시계 \\
스파이크 사건 & 발생 시각과 선택적 파형 & 전극ㆍunit 식별자, 검출기 버전 \\
\bottomrule
\end{tabularx}
\label{tab:channel_roles}
\end{table}

\subsection{표본 발생 규칙과 수치 표현}
\begin{definition}[스트림 형식]
스트림의 형식을 다음 다섯 성분으로 정의한다.
\begin{equation}
 \tau=\Stream[T_{\mathrm{val}},\mathcal C,\kappa,\mathcal R,\mathcal Q].
 \label{eq:streamtype}
\end{equation}
$T_{\mathrm{val}}$은 값의 형식, $\kappa$는 시계 영역,
$\mathcal R$은 표본 발생 규칙, $\mathcal Q$는 품질 어휘이다.
수치 기본형 $D$는 정수ㆍ부동소수점ㆍ패킹 표현을 정한다.
값 형식은 $D$, $\operatorname{Vec}[D,C]$, $\operatorname{Epoch}[D,C,N]$,
$\operatorname{SPD}[D,C]$, 사건의 합 형식으로 구성한다.
벡터와 epoch의 채널 축은 $\mathcal C$의 순서를 사용한다.
\begin{equation}
 \mathcal R=\operatorname{Uniform}(f,\phi)
 \mid\operatorname{Events},\qquad f\in\mathbb Q_{>0}.
\end{equation}
균일 스트림은 값 레코드의 명목 발생률 $f$와 위상 $\phi$를 사용한다.
epoch 값은 내부 표본 격자와 별도로 epoch 발생률을 가진다.
이벤트 스트림은 각 레코드의 시간 정보를 사용한다.
\end{definition}

획득 블록의 값 배열은 $X\in D^{C\times N}$으로 표현한다. 첫 번째 축은 채널,
두 번째 축은 표본이다. 어댑터는 입력 버퍼의 stride, endianness와 채널 순서를
이 형식에 대응시킨다. 배열의 저장 표현과 채널의 물리 단위는 각자의 기술자를 유지한다.
BrainFlow의 행렬 경로는 EEG, 보조 센서, 사건, timestamp 행을 제공하며,
각 preset이 자신의 표본률과 시간 정보를 가진다\cite{brainflow_data}.
이 경로의 어댑터는 행 역할을 분해하여 식~\eqref{eq:streamtype}의 스트림들을 생성한다.

\subsection{측정 시각, 가용 시각과 세대}
\begin{definition}[시간 레코드]
한 표본 또는 사건을 다음 튜플로 나타낸다.
\begin{equation}
 r_i=(\mathrm{id}_i,x_i,K_i,I_i^{\mathrm{meas}},a_i,
       e_i^{\mathrm{schema}},e_i^{\mathrm{clock}},q_i,\rho_i).
 \label{eq:record}
\end{equation}
$I_i^{\mathrm{meas}}$는 대표 측정 시각의 구간이고, $a_i$는 이름이 지정된 관측 경계에서
해당 레코드를 이용할 수 있게 된 시각이다. $\rho_i$는 데이터ㆍ매개변수ㆍ변환의 계보를 나타낸다.
\end{definition}

시계 변환이 점 추정 $\widehat t_i$와 오차 반경 $\epsilon_i$를 제공하면
\begin{equation}
 I_i^{\mathrm{meas}}=[\widehat t_i-\epsilon_i,\widehat t_i+\epsilon_i]
\end{equation}
로 표현한다. 획득 패킷 도착, SDK 반환, bridge 송신, DSL 수신은 각각 관측 경계 이름을 갖는다.
연산 결과는 대표 시각을 유지하면서 자신의 실제 가용 시각을 기록한다.
데이터의 전체 의존 시간 지지는 원천 조상의 측정 구간으로부터 계산한다.
\begin{equation}
 \operatorname{Supp}(y)=
 \bigcup_{r\in\operatorname{SourceAnc}(y)}I_r^{\mathrm{meas}}.
 \label{eq:support}
\end{equation}

표본 인덱스의 출처는 다음 합 형식으로 보관한다.
\begin{equation}
 K_i=\operatorname{Observed}(k_i)
 \mid\operatorname{Reconstructed}(k_i,\nu_i)
 \mid\operatorname{Unresolved}.
\end{equation}
$\nu_i$는 카운터, 표본률, 보정 구간 등 재구성에 사용한 근거이다.
수신기는 acquisition epoch에서 표본 순서의 연속성을 관리하고,
schema epoch에서 채널ㆍ스케일ㆍ표현의 연속성을 관리한다.
clock epoch는 장치 시계의 재시작과 변환 유효 구간을 식별한다.

\begin{definition}[가용 접두 스트림]
실행 환경과 안정적인 레코드 순서를 고정하고, 시각 $u$에 가용한 부분을 다음과 같이 정의한다.
\begin{equation}
 \cut_u(S)=(r_i\in S\mid a_i\le u).
 \label{eq:cut}
\end{equation}
동일한 가용 시각의 레코드는 원천 식별자와 순번으로 정렬한다.
타이머, 설정 변경, 모델 활성화도 이름이 지정된 입력 사건으로 포함한다.
\end{definition}

\subsection{품질 상태와 계보 그래프}
품질 상태는 $q_i=(v_i,o_i,F_i)$로 구성한다. $v_i$는 채널별 유효성 마스크,
$o_i$는 표본 출처, $F_i$는 사건 플래그 집합이다. 출처는
$\mathsf{Measured}$, $\mathsf{Interpolated}$, $\mathsf{Recovered}$ 등을 사용한다.
플래그는 gap, saturation, decoder resynchronization, epoch transition을 표현한다.
Cortex의 EEG 레코드는 보간 여부를 제공하고\cite{cortex_sample}, OxySoft는
연결 손실 자료의 복구 경로를 제공한다\cite{oxysoft}. 어댑터는 이 정보를
품질 상태와 가용 시각에 반영한다.

계보는 방향성 비순환 그래프로 저장한다. 원천 노드에는 장치 프로필과 원본 범위를,
변환 노드에는 연산자ㆍ매개변수ㆍ버전과 입력 노드 식별자를 기록한다.
모델 노드는 학습 자료, 전처리 통계와 선택 절차를 연결한다.
구현은 정규화된 레코드의 내용 해시를 식별자로 사용하여 반복 참조를 공유한다.
수치 연산과 메타데이터 변환이 동일한 계보 경로를 따라가므로,
출력에서 원래 측정과 활성 모델까지의 관계를 복원할 수 있다.

\section{언어 코어와 합성 의미론}
\label{sec:semantics}

언어 코어는 상태를 가진 연산자와 스트림 연결로 프로그램의 의미를 구성한다.
Kahn의 수열 의미론은 통신선의 이력과 함수 합성을 연결하며,
동기 데이터흐름과 Lustre는 표본 발생 규칙과 논리 시간에 따라 연산자 그래프를 조직한다
\citep{kahn,sdf,lustre}.
본 절은 이 구조에 측정 스키마, 가용 시각, 데이터 계보를 결합한다.
형식 판단은 한 스키마 세대의 그래프에 적용하고, 설정 사건은 새로운 스키마 세대와
그에 대응하는 상태 초기화 규칙을 선택한다.

\subsection{원시 연산자와 상태 불변식}

\begin{definition}[상태 전이 연산자]
스트림 형식 $\tau$에 대해 $\Val(\tau)$를 해당 스키마와 시간 메타데이터를
만족하는 유한 블록의 집합으로 둔다. 원시 연산자 $o$는 다음 튜플을 가진다.
\begin{equation}
 o=(\tau_o^{\mathrm{in}},\tau_o^{\mathrm{out}},Z_o,z_o^0,
       \delta_o,\mathcal J_o,\mathcal E_o).
\end{equation}
$Z_o$는 상태 공간, $z_o^0$는 초기 상태,
$\mathcal J_o\subseteq Z_o$는 상태 불변식이다. 전이 함수는 입력 블록을 소비하고
다음 상태와 유한 개의 출력 블록을 계산한다.
\begin{equation}
 \delta_o:Z_o\times\Val(\tau_o^{\mathrm{in}})
 \longrightarrow Z_o\times\Val(\tau_o^{\mathrm{out}})^*.
\end{equation}
\end{definition}

각 원시 연산자의 계약은 초기 상태의 적합성과 전이의 폐쇄성을 포함한다.
\begin{align}
 z_o^0&\in\mathcal J_o,\label{eq:primitive-init}\\
 \left.
 \begin{aligned}
 z&\in\mathcal J_o,\qquad x\in\Val(\tau_o^{\mathrm{in}}),\\
 \delta_o(z,x)&=(z',Y)
 \end{aligned}\right\}
 &\Longrightarrow
 \begin{cases}
 z'\in\mathcal J_o,\\
 Y\in\Val(\tau_o^{\mathrm{out}})^*.
 \end{cases}
 \label{eq:primitive-preservation}
\end{align}
필터 상태는 지연 표본과 내부 누산값을 보유하고, 학습 연산자는 모델 매개변수와
갱신 계보를 보유한다. 스키마에 의존하는 상태는 채널 순서, 단위, 기준,
시계 세대의 식별자를 함께 저장한다.

출력 $y$의 실제 의존 레코드 집합을 $\Dep(y)$로 둔다.
연산자는 입력, 현재 상태, 매개변수에 연결된 조상 집합을 출력과 다음 상태로 전달한다.
이 계보 계약은 다음 포함 관계를 만족한다.
\begin{equation}
 \Dep(y)\subseteq\rho(y).
 \label{eq:lineage-inclusion}
\end{equation}
조건 분기와 매개변수 선택에 사용한 자료도 이 조상 집합에 포함한다.
따라서 계보는 수치 계산과 계산 경로 선택의 의존성을 함께 나타낸다.

\subsection{효과와 코어 문법}

\begin{definition}[연산 효과]
효과 $\mathcal E_o=(m_o,R_o,h_o)$는 처리 모드
$m_o\in\{\Live,\Record\}$, 요구 기능 집합 $R_o$,
출력의 측정 레이블을 기준으로 한 전방 관측 길이 $h_o$를 기록한다.
$\Live$ 모드는 가용 입력에 따라 결과를 생성하고, $\Record$ 모드는
입력 기록의 구간과 종료 상태를 연산 계약에 포함한다.
\end{definition}

유한 전방 창을 사용하는 연산자는 창에 필요한 표본과 완결 사건을 받은 뒤 출력을 생성한다.
창의 대표 측정시각은 끝점 anchor를 사용하며, 중심 창은 중심 anchor와
전방 관측 길이 $h_o$를 함께 명시한다.
출력 레코드는 대표 측정시각, 시간 지지 $\operatorname{Supp}(y)$,
실제 가용 시각을 함께 유지한다.
형식 검사기는 연산자의 시간 지원 구간을 합성하여 프로그램의 관측 범위를 계산한다.
연산 비용, 대기 시간, 상태 크기는 실행 계획의 자원 계약으로 연결한다.

프로그램 코어의 문법은 다음과 같다.
\begin{equation}
 \begin{aligned}
 p::={}&o\mid p_1;p_2\mid p_1\otimes p_2\\
       &\mid\operatorname{delay}_{d,z}
        \mid\operatorname{feedback}(p).
 \end{aligned}
 \label{eq:core-grammar}
\end{equation}
형식 환경 $\Gamma$는 원시 연산자의 서명과 변환 계약을 보관한다.
형식 판단 $\Gamma\vdash p:\tau\xrightarrow{\mathcal E}\sigma$는
프로그램 $p$의 입력 형식, 출력 형식, 실행 효과를 연결한다.
직렬 합성은 다음 규칙을 사용한다.
\begin{equation}
 \frac{
 \Gamma\vdash p:\tau_1\xrightarrow{\mathcal E_1}\tau_2
 \qquad
 \Gamma\vdash q:\tau_2\xrightarrow{\mathcal E_2}\tau_3
 }{
 \Gamma\vdash p;q:\tau_1
 \xrightarrow{\mathcal E_1\circ\mathcal E_2}\tau_3
 }.
 \label{eq:typing-sequential}
\end{equation}
공유 형식 $\tau_2$는 채널 순서, 측정량, 단위, 시계, 발생 규칙을 일치시킨다.
단위 변환, 재기준화, 재표본화, 시계 변환은 각각의 입력과 출력을 가진
명시적 연산자로 연결한다.
직렬 의미는 함수 합성으로 정의한다.
\begin{equation}
 \den{p;q}=\den{q}\circ\den{p}.
\end{equation}

병렬 합성은 입력과 출력 형식의 곱을 사용한다.
\begin{equation}
 \frac{
 \Gamma\vdash p:\tau_1\xrightarrow{\mathcal E_1}\sigma_1
 \qquad
 \Gamma\vdash q:\tau_2\xrightarrow{\mathcal E_2}\sigma_2
 }{
 \Gamma\vdash p\otimes q:\tau_1\times\tau_2
 \xrightarrow{\mathcal E_1\otimes\mathcal E_2}
 \sigma_1\times\sigma_2
 }.
 \label{eq:typing-parallel}
\end{equation}
두 합성에서 요구 기능은 $R_1\cup R_2$로 계산한다.
구성 연산자들이 $\Live$ 모드를 가지면 합성 결과에도 $\Live$를 부여한다.
$\Record$ 연산자를 포함하는 합성에는 $\Record$를 부여한다.
가변 상태와 장치 핸들은 지정한 actor가 소유하며, 읽기 전용 블록은
수명 계약에 따라 처리 분기에 공유한다. 인터페이스별 상태와 수명 계약은
RustBelt의 의미론적 라이브러리 검증 조건과 연결된다\citep{rustbelt}.

\subsection{가용 시각 인과성과 합성 보존}

동일한 초기 상태와 실행 환경에서 프로그램의 스트림 의미를 $\den{p}$로 쓴다.
각 유한 실행 시간 구간에 속하는 입력 레코드의 수는 유한하다.
실시간 연산자는 시각 $u$까지 가용한 자료와 실행 사건으로
해당 시각까지의 출력을 결정한다.

\begin{definition}[가용 시각 인과성]
스트림 변환 $F$가 모든 입력 $S$와 실행 시각 $u$에 대해
다음 절단 법칙을 만족하면 $F$를 가용 시각에 대해 인과적이라고 정의한다.
\begin{equation}
 \cut_u(F(S))=\cut_u\bigl(F(\cut_u(S))\bigr).
 \label{eq:availability-causality}
\end{equation}
각 출력의 가용 시각은 참조한 자료의 가용 시각과 다음 관계를 가진다.
\begin{equation}
 a_y\ge\max_{r\in\Dep(y)}a_r.
 \label{eq:availability-dependencies}
\end{equation}
상수 출력은 초기 상태와 상수 매개변수의 가용 시각을 사용한다.
\end{definition}

식~\eqref{eq:availability-causality}는 입력 관측과 출력 결정을 동일한
실행 시간 절단으로 연결한다. 측정 시각을 기준으로 한 창의 레이블은
창에 속한 물리 현상을 나타내고, 가용 시각은 그 창을 이용하는 의사결정의 시점을 나타낸다.
실행 기록은 이 두 시각과 모델 활성화 사건을 함께 보존한다.

\begin{proposition}[직렬ㆍ병렬 합성의 보존]
원시 연산자들이 식~\eqref{eq:primitive-init}--\eqref{eq:primitive-preservation}의
상태 계약과 식~\eqref{eq:lineage-inclusion}의 계보 계약을 만족한다고 하자.
원시 연산자, 직렬 합성, 병렬 합성으로 구성한 형식화된 프로그램은
유효 입력을 유효 출력으로 사상하고 출력의 실제 의존성을 계보에 포함한다.
구성 원시 연산자들이 식~\eqref{eq:availability-causality}를 만족하면
합성 프로그램도 가용 시각 인과성을 만족한다.
\end{proposition}

\begin{proof}
프로그램 구조에 대한 귀납법을 적용한다. 원시 연산자 경우는 해당 계약에서 따른다.
직렬 합성에서 첫 연산자의 출력 형식과 다음 연산자의 입력 형식이 일치하므로
유효 값의 사상이 이어진다. $F=\den{p}$, $G=\den{q}$라 두면
인과성은 다음 등식으로 얻는다.
\begin{align*}
 \cut_u(G(F(S)))
 &=\cut_u\bigl(G(\cut_u(F(S)))\bigr)\\
 &=\cut_u\bigl(G(\cut_u(F(\cut_u(S))))\bigr)\\
 &=\cut_u\bigl(G(F(\cut_u(S)))\bigr).
\end{align*}
계보의 합집합과 조상 폐쇄는 두 연산 단계의 실제 의존성을 포함한다.
병렬 합성은 두 성분에 귀납 가정을 각각 적용한다.
따라서 형식 적합성, 계보 포함 관계, 가용 시각 인과성을 전체 프로그램으로 확장한다.
\end{proof}

\subsection{초기 지연과 피드백 실행}

논리 틱으로 전개한 반응형 부분그래프는 틱 $k$의 유한 입력 블록과 actor 상태로
현재 출력과 다음 상태를 계산한다. 지연 연산자는 $d\ge1$개의 틱과
형식에 적합한 초기 블록 $z_0,\ldots,z_{d-1}$를 가진다.
\begin{equation}
 \operatorname{delay}_{d,z}x[k]
 =\begin{cases}
 z_k,&0\le k<d,\\
 x[k-d],&k\ge d.
 \end{cases}
 \label{eq:initialized-delay}
\end{equation}
입출력 시계와 스키마가 연결 형식을 결정하고, 초기 블록은 그 형식을 만족한다.
양의 지연을 가진 연결선을 제거한 그래프를 $G_0$라 둔다.
피드백 형식 규칙은 $G_0$의 비순환성과 각 actor의 결정적 전이를 요구한다.
이 규칙은 현재 틱의 계산 순서와 이전 틱에서 전달할 상태를 명시한다.
형식 $\tau\times\sigma\xrightarrow{\mathcal E}\nu\times\sigma$를 가진
부분그래프의 $\sigma$ 출력과 $\sigma$ 입력을 초기 지연으로 연결하면,
외부 형식은 $\tau\xrightarrow{\mathcal E}\nu$가 된다.

\begin{proposition}[지연 피드백의 실행 유일성]
반응형 그래프의 각 actor가 유효 입력과 현재 상태에 대해
전역적으로 정의된 결정적 전이를 가진다고 하자.
초기 상태와 지연 초기값이 각 형식을 만족하고 $G_0$가 비순환 그래프이면,
각 외부 입력 수열에 대해 유일한 형식화된 실행 수열이 존재한다.
\end{proposition}

\begin{proof}
틱 $k=0$의 지연값은 초기 블록으로 정해진다.
외부 입력과 지연값을 고정하고 $G_0$의 위상 순서대로 actor를 평가하면
모든 출력과 다음 상태가 유일하게 정해진다.
틱 $k$까지의 실행이 유일하게 정해졌다고 하자.
틱 $k+1$의 지연값은 초기 블록 또는 앞선 틱의 출력에서 결정된다.
같은 위상 순서 평가가 틱 $k+1$의 출력과 상태를 유일하게 결정한다.
각 actor의 상태 계약과 연결 형식은 유효 값을 다음 틱으로 전달한다.
자연수 귀납법으로 전체 실행 수열의 존재와 유일성을 얻는다.
\end{proof}

두 명제는 처리 그래프의 합성과 폐루프 실행에 대한 수학적 기준을 제공한다.
컴파일러는 형식화된 그래프를 실행 계획으로 옮기고,
그 계획에 토큰 수, 버퍼 용량, 호출 시간과 상태 소유권을 연결한다.

\section{신호 처리 원시 연산자의 수학적 계약}
\label{sec:operators}
원시 연산자는 수치 함수와 함께 입력 스키마, 출력 스키마, 상태 불변식, 시간 지지,
초기화 규칙과 품질 전파 규칙을 제공한다. 본 절은 전기생리ㆍ광학 신호의 핵심 연산을
이 계약의 형태로 정의한다.

\subsection{단위 변환과 재기준화}
단위 변환은 채널별 선형 또는 affine 사상으로 표현한다. 전극 전압의 재기준화는
가중치 $w\in\mathbb R^C$, $w^\top\mathbf 1=1$을 사용하여 다음과 같이 계산한다.
\begin{equation}
 A_w=I-\mathbf 1w^\top,\qquad Y=A_wX.
 \label{eq:reref}
\end{equation}
연산자는 출력 기준을 $w$와 입력 채널 식별자의 조합으로 기록한다.
$A_w\mathbf 1=0$이므로 공통 성분에 해당하는 영공간을 형성한다.
후속 공분산 연산은 이 선형 변환과 자신의 정규화 계약을 함께 사용한다.
채널 선택과 순서 변경은 명시적인 선택 행렬 $P$로 표현하며 $Y=PX$와 출력 스키마 $P\mathcal C$를 생성한다.

\subsection{필터와 표본률 변환}
길이 $L$인 FIR 필터를 다음 상태 전이로 정의한다.
\begin{equation}
 y_c[n]=\sum_{\ell=0}^{L-1}h[\ell]x_c[n-\ell].
 \label{eq:fir}
\end{equation}
상태는 채널별 이전 $L-1$개 표본이다. 초기 상태의 값과 초기 유효성 구간을 계약에 기록한다.
계수 $h$가 무차원이면 출력은 입력 전압의 단위를 유지한다.
출력의 대표 시각은 $t_n$으로 정하고, 시간 지지는 $t_{n-L+1}$부터 $t_n$까지의 입력 구간으로 구성한다.
선형 위상 계수가 $h[\ell]=h[L-1-\ell]$을 만족하면 군지연은 $(L-1)/(2f_s)$이다.
실행기는 이 신호 지연과 계산 완료 시각을 각각 기록한다.

IIR 필터는 상태공간 표현을 사용한다.
\begin{equation}
 z_{n+1}=Az_n+Bx_n,\qquad y_n=Cz_n+Dx_n.
\end{equation}
필터 프로필은 계수의 수치 형식, 실현 방식, 초기 상태, 포화 정책과 안정성 조건을 가진다.
$\rho(A)<1$인 선형 시불변 실현은 유계 입력에 대한 상태의 유계성을 제공한다.
실행기는 선택한 계수와 부동소수점 실현을 버전으로 고정한다.

유리수 비율 $p/q$의 재표본화는 입력 격자와 출력 격자를 연결한다.
\begin{equation}
 f_{\mathrm{out}}=\frac pq f_{\mathrm{in}},\qquad
 y[m]=\sum_n x[n]h[mq-np],\qquad p,q\in\mathbb N_{>0}.
 \label{eq:resample}
\end{equation}
계약은 재구성 필터 $h$, 위상 원점, 경계 초기값, anti-alias 대역을 지정한다.
출력 계보는 각 표본에 기여한 입력 구간과 재표본화 설정을 연결한다.
정렬 연산자는 대상 시계를 선택하고, 재표본화 연산자는 대상 격자를 선택한다.

\subsection{창과 사건 기준 구간}
끝점 기준 창은 시작 위상 $n_0$, 길이 $N$, hop $H$, $N,H\in\mathbb N_{>0}$를 가진다.
\begin{equation}
 X_k=[x[n_0+kH],x[n_0+kH+1],\ldots,x[n_0+kH+N-1]]
      \in\mathbb R^{C\times N}.
 \label{eq:window}
\end{equation}
출력 값 형식은 $\operatorname{Epoch}[D,C,N]$이며,
출력 대표 시각은 $t_{n_0+kH+N-1}$이고, 연속 출력의 명목 간격은 $H/f_s$이다.
표본 첫 시각과 끝 시각의 차이는 $(N-1)/f_s$이다.
원시 구간의 표본 집합과 품질 마스크가 창의 유효성을 결정한다.
사건 기준 구간은 사건 $e$의 시각 $t_e$와 상대 구간 $[a,b)$를 사용하여
$[t_e+a,t_e+b)$의 표본을 선택한다. 입력 진행 경계가 구간의 끝에 도달하면 창을 완결한다.
중심 창은 대표 시각과 전방 관측 길이를 연산자 효과에 포함한다.

\subsection{전력 스펙트럼과 무차원 특징}
실수 입력과 길이 $N\ge1$의 창 함수 $w[n]$를 사용한다.
창 함수는 $\sum_{n=0}^{N-1}w[n]^2>0$을 만족한다.
이산 푸리에 변환을
\begin{equation}
 \widehat X_c[k]=\sum_{n=0}^{N-1}w[n]x_c[n]
                 e^{-2\pi\mathrm{i}kn/N}
\end{equation}
로 정의한다. 한쪽 스펙트럼의 정규화는
\begin{equation}
 P_c[k]=\frac{\eta_k|\widehat X_c[k]|^2}
                  {f_s\sum_{n=0}^{N-1}w[n]^2},\qquad
 f_k=\frac{kf_s}{N}
 \label{eq:psd}
\end{equation}
를 $k=0,\ldots,\lfloor N/2\rfloor$에 사용한다.
$\eta_k$는 DC와 짝수 $N$의 Nyquist 성분에서 1, 나머지 양의 주파수에서 2이다.
입력 단위가 $u$이면 $P_c[k]$의 단위는 $u^2/\mathrm{Hz}$이다.
주파수 대역 $B$의 전력과 로그 특징은
\begin{equation}
 E_{c,B}=\sum_{k:f_k\in B}P_c[k]\frac{f_s}{N},\qquad
 \varphi_{c,B}=\log\frac{E_{c,B}+E_0}{E_{\mathrm{ref}}}
\end{equation}
로 정의한다. 양의 상수 $E_0,E_{\mathrm{ref}}$는 $u^2$ 단위를 가지며,
특징값 $\varphi_{c,B}$는 무차원이다. 언어는 대역 경계, 창 함수, 정규화와 기준 전력을
특징 스키마에 기록한다.

\subsection{공분산과 공간 특징}
평균 제거 행렬을 $H_N=I_N-\mathbf 1\mathbf 1^\top/N$으로 두면 표본 공분산은
\begin{equation}
 \widehat\Sigma=\frac{1}{N-1}XH_NX^\top,\qquad N\ge2.
\end{equation}
공분산 연산은 동일한 물리량과 단위를 갖는 채널 묶음을 입력으로 받는다.
일반적인 다중 물리량 입력은 선언한 채널별 정규화로 무차원 벡터를 생성한 후 공분산을 계산한다.
양정치 행렬을 요구하는 후속 연산을 위해 다음 정규화를 정의한다.
\begin{equation}
 \begin{aligned}
 s^2&=\max\left(\frac{\tr\widehat\Sigma}{C},\sigma_0^2\right),
       &\sigma_0^2&>0,\\
 \Sigma_\lambda&=(1-\lambda)\widehat\Sigma+\lambda s^2I_C,
       &0&<\lambda\le1.
 \end{aligned}
 \label{eq:spd}
\end{equation}
$\sigma_0^2$는 공분산과 같은 단위를 갖는다.

\begin{proposition}[정규화 공분산의 양정치성]
$\widehat\Sigma\succeq0$, $\sigma_0^2>0$, $0<\lambda\le1$이면
$\Sigma_\lambda\succ0$이다.
\end{proposition}
\begin{proof}
임의의 $z\ne0$에 대해
\[
 z^\top\Sigma_\lambda z
 =(1-\lambda)z^\top\widehat\Sigma z+\lambda s^2\lVert z\rVert^2
 \ge\lambda\sigma_0^2\lVert z\rVert^2>0.
\]
따라서 양정치 조건이 성립한다.
\end{proof}

CSP 연산자는 두 클래스의 공분산에서 공간 가중치를 학습한다\cite{csp}.
본 언어의 CSP 계약은 정규화된 클래스 공분산 $\Sigma_1,\Sigma_2$와 다음 고유문제를 사용한다.
\begin{equation}
 \Sigma_1w_j=\lambda_j(\Sigma_1+\Sigma_2)w_j,
 \qquad w_j^\top(\Sigma_1+\Sigma_2)w_j=1.
\end{equation}
선택한 고유벡터, 입력 채널 순서와 기준을 모델 상태에 저장한다.
선택한 고유벡터의 유한 인덱스 집합을 $J$, $|J|\ge1$로 둔다.
현재 창의 정규화 공분산으로부터 다음 로그 특징을 계산한다.
\begin{equation}
 e_j=w_j^\top\Sigma_\lambda w_j,\qquad
 \varphi_j=\log\frac{e_j}{\sum_{\ell\in J}e_\ell},\qquad j\in J.
 \label{eq:csp-feature}
\end{equation}
고유벡터의 정규화와 $\Sigma_\lambda\succ0$에서 $e_j>0$을 얻는다.
분산의 비율은 무차원이며 로그 특징의 정의역을 만족한다.

Riemannian 특징은 양정치 공분산의 기하를 이용한다\cite{riemannian}.
학습 집합에서 구한 기준 $R\succ0$에 대해
\begin{equation}
 Z=\log(R^{-1/2}\Sigma_\lambda R^{-1/2}),\qquad
 \varphi=\operatorname{svec}(Z)
 \label{eq:tangent}
\end{equation}
를 사용한다. $\operatorname{svec}$는 대각 원소와 $\sqrt2$배의 상삼각 비대각 원소를
벡터로 구성하여 Frobenius 내적을 보존한다. 식~\eqref{eq:tangent}의 로그 인수는 무차원이다.
기준 $R$과 채널 스키마는 모델 계보에 속한다.

\subsection{광학 측정 변환}
광강도 $I_\lambda(t)>0$와 양의 기준 $I_{0,\lambda}$에서 광밀도 변화를
\begin{equation}
 \Delta\mathrm{OD}_\lambda(t)
 =-\log\frac{I_\lambda(t)}{I_{0,\lambda}}
\end{equation}
로 정의한다. 광경로와 흡광 계수를 포함한 행렬 $A$를 사용하면
\begin{equation}
 \Delta\mathrm{OD}=A\,\Delta c+\xi,
 \qquad
 \widehat{\Delta c}=(A^\top WA)^{-1}A^\top W\Delta\mathrm{OD}
 \label{eq:mbll}
\end{equation}
를 얻는다. 이 계약은 $W\succ0$, $\rank(A)$가 추정 농도 성분 수와 같다는 조건을 사용한다.
광학 연산자는 선택한 파장, source--detector 거리, 차등 광경로 인자,
흡광 계수의 단위와 로그 밑을 $A$의 메타데이터로 기록한다.
이 구성은 fNIRS 분석의 광밀도ㆍ농도 변환 관계를 따른다\cite{homer}.
SNIRF의 채널 대응 정보와 NIRx의 원시ㆍ파생 출력은 동일한 광학 스키마로 연결한다\cite{snirf,nirx_formats}.

\subsection{수치 실행 계약}
각 연산자는 자료형, 반올림 정책, 유효 입력 영역, 출력 오차 기준을 선언한다.
참조 연산 $F$와 실행 커널 $\widetilde F$의 비교에는
\begin{equation}
 \lVert\widetilde F(x)-F(x)\rVert
 \le\epsilon_{\mathrm{abs}}+\epsilon_{\mathrm{rel}}\lVert F(x)\rVert
 \label{eq:numeric}
\end{equation}
형태의 조건을 사용한다. 정수 디코더는 패킷 규격에 대응하는 정확한 값 관계를 사용하고,
부동소수점 커널은 선택한 입력 영역과 자료형에 맞는 수치 조건을 사용한다.
전이 상태의 유효성 검사와 품질 정책은 각 연산자의 실행 결과에 함께 기록한다.

\section{다중 하드웨어 추상화와 구성 계약}
\label{sec:hardware}
\subsection{지원 프로필과 접속 계열}
하드웨어 프로필은 실제 획득 경로의 식별자이다.
\begin{equation}
 \begin{aligned}
 H=(&\mathrm{model},\mathrm{hwrev},\mathrm{firmware},\mathrm{receiver},\\
    &\mathrm{route},\mathrm{acqVersion},\mathrm{sdkVersion},\mathrm{OS/arch}).
 \end{aligned}
 \label{eq:hardwareprofile}
\end{equation}
\code{route}는 직접 프로토콜, 직접 SDK, 획득 호스트, 기록ㆍ재생 중 하나를 선택한다.
직접 프로토콜은 장치 패킷을 해석하고, 직접 SDK는 제조사 버퍼와 제어 인터페이스를 수용한다.
획득 호스트는 RDA, RHX TCP, SpikeGLX TCP, Open Ephys ZMQ와 같은 프로그램 경계를 사용한다.
기록 어댑터는 저장 형식의 시간ㆍ표현ㆍ계보를 복원한다\cite{recorder,intan_rhx,spikeglx,open_ephys_zmq}.

Bluetooth Generic Health Sensor Service는 단위, 배율, 오프셋과 표본 배열을 표현한다.
USB Personal Healthcare Device Class는 해당 클래스의 장치 인터페이스를 정의한다\cite{ghss,usb_classes}.
프로필은 장치가 채택한 전송 프로토콜과 측정 스키마를 함께 기록한다.
Cyton의 전용 무선ㆍ시리얼 경로, Ganglion의 BLE 경로, Unicorn의 Bluetooth 2.1+EDR 경로는
각각의 접속 계약을 가진다\cite{cyton,ganglion,unicorn_manual}.

\subsection{유효 구성과 기능 집합}
프로필 $H$의 설정 공간을 $\Omega_H$, 선택 설정을 $\theta$로 둔다.
설정은 채널 집합, 획득률, 출력률, gain, reference, 필터, 입출력 모드를 포함한다.
설정별 기능 집합은 $C(H,\theta)$로 표현한다.
\begin{equation}
 \begin{aligned}
 \mathcal K=\{&\mathsf{AcquireElectrical},\mathsf{AcquireOptical},
               \mathsf{AcquireAuxiliary},\\
             &\mathsf{Configure},\mathsf{MeasureImpedance},
               \mathsf{ReceiveEvent},\mathsf{EmitEvent},\\
             &\mathsf{InsertMarker},\mathsf{MapClock},\mathsf{Stimulate}\},
 \qquad C(H,\theta)\subseteq\mathcal K.
 \end{aligned}
\end{equation}
기능의 동시 사용 관계와 정량 조건은 $K_{H,\theta}$에 기록한다.
Intan RHD 프로필은 기록ㆍ설정ㆍ보조 입출력을 기술하고, RHS 프로필은 전극 자극과
기록 상태를 함께 기술한다\cite{intan_rhd,intan_rhs}.
Neuroelectrics 경로는 EEG, 가속도, 자극 사건과 제어 명령을 기능별 인터페이스에 대응시킨다\cite{neuroelectrics}.

프로그램 $P$의 기능 요구를 $R_P$, 형식ㆍ시간ㆍ자원 요구를 $K_P$, 실행 환경 조건을 $A$로 둔다.
바인딩 판정은 다음 결합 조건을 평가한다.
\begin{equation}
 \begin{aligned}
 \operatorname{Bind}(P,H,\theta,A)\iff{}&
 \theta\in\Omega_H\ \land\ R_P\subseteq C(H,\theta)\\
 &\land\operatorname{SAT}(K_{H,\theta}\land A)\\
 &\land(K_{H,\theta}\land A\models K_P).
 \end{aligned}
 \label{eq:bind}
\end{equation}
판정기는 각 절의 결과를 $\mathsf{Pass}$, $\mathsf{Fail}$, $\mathsf{Pending}$으로 기록한다.
$\mathsf{Pending}$ 항목은 필요한 구성값과 증거 식별자를 수집하고,
$\mathsf{Fail}$ 항목은 해당 입력ㆍ출력 계약을 진단 정보로 반환한다.
이 구조는 문서 제약, 런타임 조회와 실행 증거를 하나의 판정식에 연결한다.

\subsection{규격에서 도출한 구성 사례}
actiCHamp의 채널ㆍ최대 표본률 조합은 표~\ref{tab:actichamp}과 같다\cite{actichamp}.
프로필은 이 상한과 장치가 제공하는 이산 sampling mode를 함께 저장한다.
\begin{table}[tb]
\centering\small
\caption{actiCHamp의 EEG 채널 구성별 최대 표본률. 각 구성은 AUX 8개를 함께 가진다.}
\begin{tabular}{@{}lr@{}}
\toprule EEG 채널 수 & 최대 표본률 \\\midrule
32 & $100\,\mathrm{kHz}$ \\
64 & $50\,\mathrm{kHz}$ \\
96, 128, 160 & $25\,\mathrm{kHz}$ \\
\bottomrule
\end{tabular}
\label{tab:actichamp}
\end{table}

Cyton+Daisy의 기본 무선 경로는 두 bank를 교대로 전송하며 연속 두 표본을 평균한다\cite{cyton}.
bank $b$의 출력값과 표본률 관계는 다음과 같이 기록한다.
\begin{equation}
 y_b[k]=\frac{x_b[n_b(k)-1]+x_b[n_b(k)]}{2},
 \quad f_{\mathrm{ADC}}=250\,\mathrm{Hz},
 \quad f_{\mathrm{out},b}=125\,\mathrm{Hz}.
\end{equation}
$n_b(k)$는 bank의 교대 위상을 포함한다. 프로필은 평균 연산과 출력 대표 시각을
별도 필드로 보관한다.

Ganglion의 firmware 3.0.0은 delta 압축에서 MSB 전송으로 코덱을 변경하였다\cite{ganglion_release}.
프로필의 스키마 식별자는 모델, firmware tag, codec과 packet layout의 조합으로 구성한다.
BioSemi ActiveThree는 A3 전용 ring buffer 형식과 $16\,384\,\mathrm{Hz}$의 장치 전송률을
사용한다\cite{biosemi_a3,biosemi_rate}. 어댑터는 세대별 디코더와 출력 decimation 설정을
각각 선택하고, 설정 변경 사건에 새 schema epoch를 부여한다.

Intan 파일 형식은 전극 값의 정수 표현을 다음 변환에 연결한다\cite{intan_file}.
\begin{equation}
 \begin{aligned}
 \texttt{.rhd}:\quad &V_{\mu\mathrm V}=0.195(r_{\mathrm{uint16}}-32768),\\
 \texttt{amplifier.dat}:\quad &V_{\mu\mathrm V}=0.195r_{\mathrm{int16}}.
 \end{aligned}
\end{equation}
디코더는 입력 형식 식별자, 부호 해석과 단위 변환을 함께 선택한다.
변환 결과와 원래 표현은 같은 계보 노드를 공유한다.

\subsection{프로필 메타데이터와 접근 조건}
BrainFlow 어댑터는 native core, Rust binding, board ID, preset과 채널 기술자를 기록한다.
Muse 2 경로는 기본 EEG 4채널과 IMU를 제공하고, 선택한 설정에서 PPG를 활성화한다.
어댑터는 각 preset의 표본률과 시계 정보를 독립 기술자로 구성한다\cite{brainflow_boards}.
LSL 어댑터는 source 식별자, outlet 인스턴스, 데이터 형식, 출력 메타데이터와
clock mapping 버전을 기록한다\cite{brainflow_build,lsl_bindings,lsl_time}.
CereLink 경로는 standalone, CereLink client, Central client의 연결 모드를 프로필에 둔다\cite{cerelink}.

접근 조건은 소프트웨어 계약과 연결한다. Brain Products SDK 프로필은 사용 조건과
상업 배포 계약을 식별하고\cite{brainproducts_sdk}, Cortex 프로필은 EEG scope와
세션 활성화 상태를 기록한다\cite{cortex_subscription}.
기능 조회와 접근 조건을 함께 평가하면 실행기가 현재 세션에서 사용할 인터페이스를 결정한다.
부록~\ref{app:hardware}는 18개 대표 모델ㆍ획득 경로를 이 프로필 구조에 대응시킨다.

\section{시계 변환과 시간 완결 계약}
\label{sec:time}

시간 계층은 각 원천의 대표 측정시각을 기준 시계로 옮기고, 가용 시각에 따라
창의 완결을 판정한다. LSL은 송신 호스트 시계와 호스트 사이의 offset 관측을 제공하며,
BrainFlow는 장치별 timestamp 의미와 preset별 표본률을 제공한다
\citep{lsl_time,brainflow_data}. 어댑터는 이 관측의 anchor, 시계 영역과 변환 버전을
시간 레코드에 연결한다.

\subsection{유효 구간을 가진 시계 사상}

원천 시계 $\kappa$의 읽기값을 $u$, 기준 시계를 $\kappa_0$로 둔다.
시계 세대 안의 변환 버전 $v$는 유효 구간 $J_v$와 중심 $u_0$를 가진다.
그 구간에서 실제 사상과 추정 사상을 다음과 같이 표현한다.
\begin{equation}
 \begin{aligned}
 \phi_v(u)&=b_v+a_v(u-u_0)+r_v(u),\\
 \widehat\phi_v(u)&=\widehat b_v+\widehat a_v(u-u_0),\qquad u\in J_v.
 \end{aligned}
 \label{eq:clock-affine}
\end{equation}
계약은 실제 사상의 증가성, $\widehat a_v>\delta_a\ge0$과 계수ㆍ잔차의 상한
$|b_v-\widehat b_v|\le\delta_b$,
$|a_v-\widehat a_v|\le\delta_a$,
$|r_v(u)|\le\delta_r$를 포함한다. 따라서 시간 구간은
\begin{equation}
 \begin{aligned}
 \epsilon_v(u)&=\delta_b+|u-u_0|\delta_a+\delta_r,\\
 \Phi_v(u)&=[\widehat\phi_v(u)-\epsilon_v(u),
                  \widehat\phi_v(u)+\epsilon_v(u)]
 \end{aligned}
 \label{eq:clock-interval}
\end{equation}
로 계산한다. 원천 timestamp의 양자화와 jitter를 구간 $U$로 표현하면,
다음 구간 확장을 사용한다.
\begin{equation}
 \Phi_v(U)=\operatorname{hull}\!\bigcup_{u\in U}\Phi_v(u).
 \label{eq:clock-interval-extension}
\end{equation}
변환 레코드는 계수, $J_v$, 가용 시각, clock epoch와 근거를 함께 저장한다.
근거 형식 $\mathsf{Bound}(A,J_v,\epsilon_v)$는 가정 $A$ 아래의 포함 관계를,
$\mathsf{Estimate}(\mathcal D,p,\widehat q_p,n)$는 자료 $\mathcal D$에서 얻은
$p$-분위수 추정과 표본 수를 기술한다.

동일한 affine 구간의 두 읽기값 $u_1,u_2$는 공통 offset을 공유하므로,
상대 시간의 오차는 다음 상한을 갖는다.
\begin{equation}
 \left|
 \phi_v(u_2)-\phi_v(u_1)-\widehat a_v(u_2-u_1)
 \right|
 \le\delta_a|u_2-u_1|+2\delta_r.
 \label{eq:clock-relative}
\end{equation}
두 사상 $\Phi_{12},\Phi_{23}$을 합성할 때 $I_{12}(u)=\Phi_{12}(u)
\subseteq J_{23}$를 요구한다. 중간 구간 전체에 걸친 오차 전파는
\begin{equation}
 \epsilon_{13}(u)\le
 |\widehat a_{23}|\,\epsilon_{12}(u)
 +\sup_{z\in I_{12}(u)}\epsilon_{23}(z)
 \label{eq:clock-compose}
\end{equation}
로 계산한다. 이 규칙은 장치, SDK 호스트와 실행기 시계 사이의 변환 사슬에 적용한다.
Rust의 \code{Instant}는 호스트에서 두 관측 사이의 경과시간을 표현한다.
실행기는 origin과 clock-domain 식별자를 부여하고 외부 시계와의 대응을
식~\eqref{eq:clock-affine}의 레코드로 관리한다\citep{rust_instant}.

\subsection{표본 격자, 패킷과 시간 anchor}

균일 발생 규칙은 원천 시계의 격자
$u_n=u_0+(n-n_0)\Delta$를 사용한다.
affine 구간의 명목 기준 표본률은
$\widehat f_{\mathrm{ref}}=(\widehat a_v\Delta)^{-1}$이다.
스트림 기술자는 채널별 ADC 획득률 $f_{\mathrm{ADC}}$,
장치 방출률 $f_{\mathrm{emit}}$, 어댑터 출력률 $f_{\mathrm{out}}$을 각각 기록한다.
패킷당 채널별 표본 수 $M$이 일정한 경로는
\begin{equation}
 f_{\mathrm{packet}}=\frac{f_{\mathrm{emit}}}{M}
 \label{eq:packet-rate}
\end{equation}
을 만족한다. 패킷 sequence와 식~\eqref{eq:record}의 표본 인덱스 $K_i$는
각자의 출처를 유지한다. 디지털 필터와 재표본화는 출력 격자,
대표시각 anchor와 의존 시간 지지를 함께 갱신한다.

장치의 ADC 관측, SDK 수신, 소프트웨어 명령과 물리 edge는 각각 이름을 가진 anchor이다.
예를 들어 PC 수신 시각을 제공하는 BrainFlow 경로는 해당 수신 경계를
timestamp anchor로 기록한다\citep{brainflow_data}.
동일 pulse의 관측 edge를 대응시키는 경로는 edge 쌍으로 clock mapping을 구성한다.
SpikeGLX의 동기화 경로는 stream별 clock calibration과 공통 sync pulse의
edge 대응을 사용한다\citep{spikeglx_manual}.

\subsection{관측 오차와 대기 자료를 포함한 watermark}

입력 $c$의 원시 레코드에 대해
$I_i^{\mathrm{meas}}=[\ell_i,h_i]
=[\widehat t_i-\epsilon_i,\widehat t_i+\epsilon_i]$로 둔다.
ingress 시각의 관측 구간을
$[\widehat a_i-\eta_i,\widehat a_i+\eta_i]$로 둔다.
한 clock epoch와 시간 뷰에서 지연 계약을 다음과 같이 정의한다.
\begin{equation}
 (\widehat a_i+\eta_i)-(\widehat t_i-\epsilon_i)\le L_c.
 \label{eq:interval-lateness}
\end{equation}
관측 중심시각의 차이가 $\widehat a_i-\widehat t_i\le L_c^0$이고
$\eta_i\le\overline\eta_c$, $\epsilon_i\le\overline\epsilon_c$이면,
$L_c=L_c^0+\overline\eta_c+\overline\epsilon_c$를 사용한다.
따라서 $L_c$는 두 관측의 시계 오차를 함께 포함한다.

현재 기준 시각의 관측 구간을 $s=[s^-,s^+]$로 둔다.
ingress 이후 join 상태에 반영되는 과정에 있는 레코드 집합을
$\operatorname{Pending}_c(s)$라 두고 다음 진행 경계를 계산한다.
\begin{equation}
 \begin{aligned}
 w_c^{\mathrm{in}}(s)&=s^--L_c,\\
 p_c(s)&=\inf\{\ell_i:i\in\operatorname{Pending}_c(s)\},
                  \qquad\inf\varnothing=+\infty,\\
 w_{\Join}(s)&=\min_c\min\{w_c^{\mathrm{in}}(s),p_c(s)\}.
 \end{aligned}
 \label{eq:pending-watermark}
\end{equation}
현재 시각 이후 도착하는 자료는 식~\eqref{eq:interval-lateness}에 따라
$\ell_i>s^--L_c$를 만족한다. 이미 도착한 자료는 $p_c(s)$가 포괄한다.
그러므로 $w_{\Join}$는 join에 앞으로 반영될 자료의 대표시각 하한이다.
구현은 자료, gap, watermark와 epoch 종료 사건의 전달 순서에 이 관계를 유지한다.

창 $W_k=[u_k,v_k)$는 다음 조건에서 완결된다.
\begin{equation}
 \begin{aligned}
 \operatorname{Complete}(W_k,s)&\iff w_{\Join}(s)\ge v_k,\\
 \operatorname{Definite}(i,W_k)&\iff I_i^{\mathrm{meas}}\subseteq W_k,\\
 \operatorname{Possible}(i,W_k)&\iff
                         I_i^{\mathrm{meas}}\cap W_k\ne\varnothing.
 \end{aligned}
 \label{eq:window-complete}
\end{equation}
연산자는 확정ㆍ가능 소속 중 선택한 정책과 경계 모호성 상태를 보관한다.
원시 레코드의 창 할당은 대표 측정 구간을 사용하고,
출력의 전체 시간 지지는 식~\eqref{eq:support}의 조상 합집합을 사용한다.
시간구간 창 출력은 오른쪽 경계 $v_k$를 대표시각 anchor로 선택한다.
표본수 창은 식~\eqref{eq:window}의 마지막 표본 시각을 anchor로 사용한다.
중심 anchor를 사용하는 연산자는 전방 관측 길이 $h_o$를 효과에 기록한다.

ingress 이후 join 반영까지의 상한이 $D_c$이고,
진행 경계 갱신 간격의 상한이 $G$이면 창의 가용 시각은
\begin{equation}
 a_{\mathrm{close}}\le v_k+\max_c(L_c+D_c)+G
 \label{eq:window-availability-bound}
\end{equation}
를 만족한다. 이 식은 기준 시각을 직접 읽는 진행 타이머를 가정한다.
구간 타이머를 사용하는 실행기는 현재 시각의 하한과 실제 시각 사이의 상한을 $G$에 더한다.
가용 시각 인과성은 창 완결 사건을 입력으로 포함하여
식~\eqref{eq:availability-causality}와 연결한다.

\subsection{Marker와 반응시간 계약}

marker는 사건 코드, 발생 원천, 관측 anchor, edge 극성,
측정 구간, 가용 시각과 clock epoch를 가진다.
장치의 디지털 입력 경로는 전압 임계값, 최소 pulse 폭과 표본화 규칙을
하드웨어 계약에 연결한다. 소프트웨어 명령 구간 $I^{\mathrm{cmd}}$와
명령에서 물리 edge까지의 지연 구간 $\Delta_{\mathrm{path}}=[d^-,d^+]$가 주어지면,
명령에 근거한 edge 예측 구간을
\begin{equation}
 I^{\mathrm{edge,pred}}=I^{\mathrm{cmd}}\oplus\Delta_{\mathrm{path}},
 \qquad
 [a,b]\oplus[c,d]=[a+c,b+d]
 \label{eq:marker-delay}
\end{equation}
로 정의한다. edge를 획득한 표본은 그 표본 격자와 clock mapping으로
자신의 $I^{\mathrm{meas}}$를 얻는다.
기준 marker 구간 $[\ell_m,h_m]$, 출력 가용 시각 상한 $\overline a_y$,
반응시간 요구 $D$에 대해
$\overline a_y\le\ell_m+D$를 충분조건으로 사용한다.
실행기는 적용한 지연 구간, 시계 변환과 창 소속 정책을 결과의 계보에 기록한다.

\section{학습, 평가와 모델 활성화}
\label{sec:learning}
\subsection{분할 단위와 데이터 역할}
학습 프로그램은 데이터의 역할과 분할 단위를 선언한다. 역할 집합을
$\{\mathsf{Fit},\mathsf{Tune},\mathsf{Eval},\mathsf{Adapt}\}$로 두고,
각 레코드에 피험자, 세션, 시행, 원본 구간 식별자를 부여한다.
분할 함수 $g$는 프로토콜에서 선택한 독립 단위를 반환한다.
\begin{equation}
 g:\mathrm{RecordID}\longrightarrow
 \mathrm{SubjectKey}\ \mid\ \mathrm{SessionKey}\ \mid\ \mathrm{TrialKey}.
\end{equation}
동일 원본 구간에서 생성한 창들은 공통 원천 식별자를 계보에 보관한다.
세션을 가로지르는 일반화 평가는 세션 key를,
피험자를 가로지르는 평가는 피험자 key를 분할 단위로 사용한다.
MOABB의 데이터셋ㆍ전처리ㆍ평가 구성을 이 데이터 역할의 연결 지점으로 사용한다\cite{moabb}.

\subsection{모델 객체와 fit 효과}
모델 버전 $v$를 다음 객체로 정의한다.
\begin{equation}
 M_v=(\vartheta_v,\mathcal C_{\mathrm{in}},\mathcal C_{\mathrm{out}},
      \rho_v^{\mathrm{fit}},a_v,\mathrm{digest}_v).
 \label{eq:model}
\end{equation}
$\vartheta_v$는 수치 매개변수, $\rho_v^{\mathrm{fit}}$는 매개변수와 선택 절차에
영향을 준 자료의 조상 집합, $a_v$는 해당 모델을 사용할 수 있게 된 시각이다.
공분산 기준, 정규화 통계, 특징 선택, 하이퍼파라미터 선택도 이 조상 집합에 포함한다.

학습 효과는 입력 자료의 역할을 소비하고 모델의 학습 계보를 생성한다.
고정 모델 평가의 적합 조건은 다음과 같다.
\begin{equation}
 \begin{aligned}
 \rho_v^{\mathrm{fit}}&\subseteq
 \Anc(D_{\mathrm{fit}}\cup D_{\mathrm{tune}}),\\
 g(\rho_v^{\mathrm{fit}})\cap g(D_{\mathrm{eval}})&=\varnothing.
 \end{aligned}
 \label{eq:split}
\end{equation}
컴파일러는 fit 효과가 생성한 모델 식별자를 추적하고, 평가 연산자는
식~\eqref{eq:split}과 입력 스키마를 검사한다.
이 조건은 필터ㆍ특징ㆍ추정기의 선택을 포함한 전체 학습 그래프에 적용한다.

중첩 평가에서는 외부 fold $j$마다 모델 $M^{(j)}$를 생성한다.
내부 fitㆍtune 자료는 외부 학습 집합의 부분집합으로 구성하고,
외부 평가 집합은 해당 fold의 결과 계산에 사용한다.
\begin{equation}
 \begin{aligned}
 D_{\mathrm{fit}}^{(j,k)}\cup D_{\mathrm{tune}}^{(j,k)}
 &\subseteq D_{\mathrm{train}}^{(j)},\\
 g(D_{\mathrm{train}}^{(j)})\cap g(D_{\mathrm{eval}}^{(j)})
 &=\varnothing.
 \end{aligned}
\end{equation}
분할 manifest는 원천 식별자, 그룹 함수, seed, 선택한 구간과 생성 절차의 버전을 보관한다.

\subsection{온라인 적응의 시간 의미}
적응 모델은 갱신 사건과 활성화 사건을 가진다. 버전 $v$의 활성화 시각을 $b_v$로 두면
\begin{equation}
 \max_{r\in\rho_v^{\mathrm{fit}}}a_r\le a_v\le b_v.
 \label{eq:modeltime}
\end{equation}
예측 작업은 시작 시점에 활성 모델 $M_{v(n)}$를 선택하고,
그 버전의 매개변수를 작업 완료까지 읽기 전용으로 보유한다.
\begin{equation}
 \widehat y_n=f_{\vartheta_{v(n)}}(x_n),\qquad
 b_{v(n)}\le a_n^{\mathrm{start}},\qquad
 \rho(\widehat y_n)=\rho(x_n)\cup\rho_{v(n)}^{\mathrm{fit}}.
\end{equation}
레이블을 사용하는 갱신은 레이블의 가용 시각을 추가 의존성으로 포함한다.
prequential 프로토콜은 예측, 결과 평가, 레이블 기반 갱신의 사건 순서를 기록한다.
비지도 적응은 입력 특징과 갱신 정책을 자신의 fit 계보에 기록한다.
정규화, 시계 보정과 calibration도 버전과 활성화 시각을 가지므로 같은 실행 규칙을 사용한다.

\subsection{추정값과 제어 사건}
추정기 출력은 class score, 확률 벡터, 연속 상태 추정량 등의 명목 형식을 갖는다.
확률 벡터는 $p_k\ge0$, $\sum_kp_k=1$을 만족하며,
확률 보정 연산은 방법ㆍ자료ㆍ버전을 추가한다.
의사결정 연산은 입력 형식과 임계값 정책을 사용하여 사건을 생성한다.
예를 들어
\begin{equation}
 d(p)=
 \begin{cases}
 \operatorname{Select}(\arg\max_kp_k),&\max_kp_k\ge\eta,\\
 \operatorname{Hold},&\max_kp_k<\eta
 \end{cases}
\end{equation}
로 정의한다. 최댓값의 동점은 선언한 클래스 순서의 첫 항목으로 결정한다.
임계값 $\eta$의 선택 자료는 해당 정책의 학습 계보에 포함한다.
명령 출력은 사건의 단위ㆍ범위ㆍ목표 시각과 장치 기능 요구를 입력받는다.
실행기는 출력한 명령과 장치 응답을 동일한 사건 식별자로 연결한다.

\subsection{평가 결과의 구조}
평가 연산자는 분할 manifest, 모델 버전, 예측ㆍ레이블 쌍, 품질 선택 정책과 지표 정의를 반환한다.
클래스별 recall을 평균하는 balanced accuracy는
\begin{equation}
 \mathrm{BA}=\frac1K\sum_{k=1}^{K}
 \frac{\mathrm{TP}_k}{\mathrm{TP}_k+\mathrm{FN}_k}
\end{equation}
로 계산하며, 각 분모가 양수인 평가 구성을 사용한다.
집계는 시행, 세션, 피험자 중 지정한 단위를 따라 수행한다.
구간 추정은 재표집 단위, 표본 수, 신뢰수준과 난수 seed를 결과 메타데이터로 포함한다.
지연 결과는 입력 준비, 출력 완료와 사건 응답의 관측 경계를 구분하여 기록한다.

\section{Rust 실행 구조와 유한 자원 계획}
\label{sec:runtime}

\subsection{형식화된 중간 표현과 소유권}

컴파일러는 구문과 이름 해석으로 연산자 그래프를 만들고,
형식 판단 $\Gamma\vdash p:\tau\xrightarrow{\mathcal E}\sigma$를 각 연결에 부여한다.
형식화된 중간 표현은 채널 스키마, 시계 사상, 발생 규칙,
품질 정책과 데이터 계보를 보유한다. 하드웨어 연결은 유효 설정과 capability 조건을,
실행 계획은 버퍼ㆍ호출시간ㆍ자원 배분 조건을 갖는다.
\begin{equation}
 \begin{aligned}
 \mathrm{Source}&\longrightarrow\mathrm{AST}
 \longrightarrow\mathrm{TypedGraph},\\
 \mathrm{TypedGraph}&\longrightarrow
        \mathrm{BoundGraph}(H,\theta)
 \longrightarrow\mathrm{ExecutionPlan}.
 \end{aligned}
 \label{eq:compiler-pipeline}
\end{equation}
각 산출물은 앞 단계의 노드 식별자와 계약을 유지한다.
격자ㆍ단위ㆍ시계 변환은 명시적 노드로 존재하며,
상태 초기화와 모델 활성화 사건은 실행 의존성에 포함된다.
코드 생성은 이 계획의 actor, 버퍼와 전이 관계를 Rust 구성 요소에 대응시킨다.

장치 actor는 native handle, SDK callback 등록, 설정과 획득 세대를 소유한다.
공통 수명 상태는 다음 전이로 구성한다.
\begin{equation}
 \begin{aligned}
 \mathsf{Created}&\longrightarrow\mathsf{Prepared}(\theta)
 \longrightarrow\mathsf{Streaming}(e),\\
 \mathsf{Streaming}(e)&\longrightarrow\mathsf{Draining}(e)
 \longrightarrow\mathsf{Closed}.
 \end{aligned}
 \label{eq:device-lifecycle}
\end{equation}
각 전이는 SDK의 완료 응답과 자원 소유 조건을 사용한다.
종료 전이는 callback 정리와 진행 중 호출의 완료를 포함한다.
설정 변경은 새 계약과 세대를 선택하는 제어 사건으로 작동한다.
fault 사건은 원인과 마지막 확정 상태를 보관하고,
복구 전이는 선택한 초기화 계약으로 새 세대를 생성한다.

가변 actor 상태는 한 소유자가 갱신하고,
완성된 데이터 블록은 소유권 이동 또는 읽기 공유로 전달한다.
\code{Arc<T>}를 사용하는 공유 객체는 $T$의 스레드 전달ㆍ공유 조건을 충족한다
\citep{rust_arc}. FFI wrapper는 포인터 유효성, 정렬, 길이, 수명과 aliasing 조건을
검증하여 native 버퍼를 형식화된 블록으로 연결한다\citep{rust_reference}.
이 경계별 검증 조건은 RustBelt가 제시한 의미론적 라이브러리 확장의
검증 구조를 따른다\citep{rustbelt}.

\subsection{고정 발생률 부분그래프의 실행 증명서}

한 설정과 스키마 세대에서 생산ㆍ소비 token 수가 일정한 부분그래프를
동기 데이터흐름 그래프 $G=(V,E)$로 나타낸다\citep{sdf}.
간선 $e$의 생산 수를 $p_e>0$, 소비 수를 $c_e>0$,
초기 token 수를 $d_e\ge0$으로 둔다.
양의 정수 반복벡터 $\boldsymbol\nu$는 균형식을 만족한다.
\begin{equation}
 p_e\nu_{\mathrm{src}(e)}
 =
 c_e\nu_{\mathrm{dst}(e)},
 \qquad e\in E.
 \label{eq:sdf-balance}
\end{equation}
순차 발화열 $\omega=(v_1,\ldots,v_M)$에서 $N_v(j)$는
첫 $j$회 발화 중 actor $v$의 횟수이다.
한 발화는 입력을 소비하고 상태 전이를 완료한 뒤 출력을 생산한다.
접두 발화열 이후의 점유량은
\begin{equation}
 m_e(j)=d_e+p_eN_{\mathrm{src}(e)}(j)
                 -c_eN_{\mathrm{dst}(e)}(j)
 \label{eq:sdf-occupancy}
\end{equation}
이다. 실행 증명서는 다음 조건을 모두 포함한다.
\begin{equation}
 \begin{aligned}
 N_v(M)&=\nu_v,\qquad v\in V,\\
 m_e(j-1)&\ge c_e,
       \qquad e:\mathrm{dst}(e)=v_j,\\
 B_e&\ge\max_{0\le j\le M}m_e(j).
 \end{aligned}
 \label{eq:sdf-certificate}
\end{equation}
입력 token 조건은 각 발화의 실행 가능성을 제공하고,
균형식은 $m_e(M)=d_e$를 준다.
각 actor의 전이가 유한 시간에 완료되고 외부 입력이 지정한 가용 조건을 만족하면,
실행기는 같은 발화열을 반복하며 용량 $B_e$ 안에서 흐름을 유지한다.
병렬 계획은 발화 시작 시 입력을 차감하고 완료 시 출력을 더하는 사건열을 사용하며,
허용한 사건열의 접두사마다 같은 token 조건과 용량 조건을 검사한다.

간선별 token 크기 상한이 $s_e$ byte이고 독립 저장 영역을 예약하면,
전체 메모리 예약량은 다음 식으로 계산한다.
\begin{equation}
 M_{\mathrm{reserve}}
 =
 \sum_{e\in E}B_es_e+
 \sum_{v\in V}M_v^{\mathrm{work}}+
 M^{\mathrm{meta}}.
 \label{eq:memory-reservation}
\end{equation}
$M_v^{\mathrm{work}}$는 actor가 동시에 소유하는 입력ㆍ출력ㆍ작업 영역의 상한이며,
$M^{\mathrm{meta}}$는 큐 관리와 제어 사건의 저장량이다.
공유 블록을 사용하는 계획은 공유 풀의 수명과 점유량으로 해당 항을 계산한다.
이벤트 marker와 설정 제어는 사건 actor로 처리하고,
고정률 영역의 경계에서 가용 조건과 세대 전이를 연결한다.

\subsection{호출시간과 주기 실행 조건}

순차 계획에 전용 실행 구간을 배정하고 actor $v$의 호출시간 상한을 $C_v$로 둔다.
발화 $j$의 의존 자료 가용 시각 상한을 $\overline a_j^{\mathrm{dep}}$,
호출ㆍ스케줄링 부가시간 상한을 $o_j$로 둔다.
주기 $k$의 시작이 $kP$일 때 완료 시각의 상한을 재귀적으로 계산한다.
\begin{equation}
 \begin{aligned}
 \overline f_{k,0}&=kP,\\
 \overline f_{k,j}
 &=\max\{\overline f_{k,j-1},
                 \overline a_{k,j}^{\mathrm{dep}}\}
                 +C_{v_j}+o_j.
 \end{aligned}
 \label{eq:cyclic-response}
\end{equation}
각 발화의 상대 deadline을 $D_j$로 두면,
\begin{equation}
 \overline f_{k,j}\le kP+D_j,
 \qquad
 \overline f_{k,M}\le(k+1)P
 \label{eq:cyclic-deadline}
\end{equation}
가 해당 순차 계획의 충분조건이다.
주기 시작 시 모든 외부 의존 자료가 준비되는 경우에는
$\sum_v\nu_vC_v+\sum_jo_j\le P$로 주기 조건을 계산한다.
실행 프로필은 $C_v$, $o_j$와 가용 상한의 근거,
CPU 배분과 OS 조건을 함께 보관한다.

\subsection{도착ㆍ서비스 곡선의 버퍼와 지연 상한}

전송 또는 처리 경계의 누적 유입량을 $A(t)$, 누적 완료량을 $D(t)$로 둔다.
두 양은 같은 payload 계량과 기준 시계를 사용한다.
초기 상태가 빈 FIFO에서 도착 곡선 $\alpha$와 서비스 곡선 $\beta$를 다음과 같이 정의한다.
\begin{equation}
 \begin{aligned}
 A(t)-A(s)&\le\alpha(t-s),\qquad 0\le s\le t,\\
 D(t)&\ge(A\otimes\beta)(t),\\
 (A\otimes\beta)(t)&=
       \inf_{0\le s\le t}\{A(s)+\beta(t-s)\}.
 \end{aligned}
 \label{eq:network-curves}
\end{equation}
burst 상한 $b$, 지속 유입률 상한 $r$,
해당 흐름에 배정한 유효 서비스율 $R>r$와 지연 $T$를 사용한다.
\begin{equation}
 \begin{aligned}
 \alpha(0)&=0,\qquad \alpha(u)=b+ru\quad(u>0),\\
 \beta(u)&=R(u-T)^+.
 \end{aligned}
 \label{eq:rate-latency}
\end{equation}
$r$는 기준 시계에서의 유입 상한이며, $b$는 패킷화와 burst를 포함한다.
$R,T$는 선택한 완료 경계까지의 자원 배분과 경합을 반영한다.
완성 블록을 계수하는 경계에서는 $T$에 블록 완성과 packetization의 지연을 포함한다.
Network calculus의 backlogㆍdelay 정리에 따라\citep{networkcalculus}
\begin{equation}
 \begin{aligned}
 A(t)-D(t)&\le b+rT,\\
 W(t)&=\inf\{d\ge0:A(t)\le D(t+d)\}
          \le T+\frac bR
 \end{aligned}
 \label{eq:network-bounds}
\end{equation}
를 얻는다. 동일한 payload 계량을 보존하는 직렬 경로는
$\beta_{\mathrm{path}}=\beta_1\otimes\cdots\otimes\beta_m$으로 합성하며,
rate--latency 모형에서는
$R_{\mathrm{path}}=\min_jR_j$,
$T_{\mathrm{path}}=\sum_jT_j$를 사용한다.
실행 계획은 이 상한을 ingress 이후의 가용 지연과 연결하여
식~\eqref{eq:window-availability-bound} 및 marker의 반응시간 조건을 계산한다.

\section{설계 식의 적용 사례}
\label{sec:cases}
\subsection{EEG 창 처리와 모델 입력}
8개 전압 채널을 $250\,\mathrm{Hz}$로 획득하는 프로필을 선택한다.
Cyton의 기본 무선 경로가 이 채널ㆍ표본률 조합을 제공한다\cite{cyton}.
분석 그래프는 전압 단위 변환, 길이 63의 대칭 FIR, 길이 250ㆍhop 25의 창,
정규화 공분산, 접공간 특징, 고정 모델 추정으로 구성한다.
FIR 계수와 학습된 접공간 기준은 그래프 매개변수 식별자로 고정한다.

창의 갱신 간격과 첫 표본부터 마지막 표본까지의 시간은
\begin{equation}
 T_{\mathrm{hop}}=\frac{25}{250}=0.1\,\mathrm{s},\qquad
 T_{\mathrm{span}}=\frac{249}{250}=0.996\,\mathrm{s}
\end{equation}
이다. FIR의 신호 군지연은 $31/250=0.124\,\mathrm{s}$이다.
초기 유효성 정책은 62개 이전 표본을 채운 후 필터 출력을 유효로 판정한다.
창의 시작 위상을 $n_0=62$로 선택하면,
완전히 유효한 첫 창의 끝은 원천 index $62+249=311$에 위치한다.
첫 원천 표본의 시각을 0으로 두면 해당 끝점은 $1.244\,\mathrm{s}$이다.

이 그래프의 형식 연결은 표~\ref{tab:eegtypes}와 같다.
\begin{table}[tb]
\centering\small
\caption{EEG 사례의 형식과 갱신 규칙}
\begin{tabularx}{\linewidth}{@{}P{.20\linewidth}YY@{}}
\toprule 연산 결과 & 값의 형식 & 시간ㆍ상태 계약 \\\midrule
전압 표본 & $\mathbb R^{8}$, V & $250\,\mathrm{Hz}$, 시계 $\kappa$ \\
FIR 출력 & $\mathbb R^{8}$, V & 채널별 이전 62표본 상태 \\
창 & $\mathbb R^{8\times250}$, V & hop 25, 끝점 anchor \\
공분산 & $\mathrm{SPD}(8)$, $\mathrm V^2$ & 창당 1개, 정규화 식~\eqref{eq:spd} \\
접공간 특징 & $\mathbb R^{36}$, 무차원 & 기준 $R$과 모델 입력 스키마 \\
예측 & 선언한 class score 또는 확률 & 활성 모델 버전과 가용 시각 \\
\bottomrule
\end{tabularx}
\label{tab:eegtypes}
\end{table}

수신 블록의 마지막 표본을 기준으로 입구 지연 4 ms, 큐 대기 2 ms,
연산 3 ms, 출력 전달 1 ms를 실행 계약으로 선택하면 끝점 이후의 반응시간은
\begin{equation}
 D_{\mathrm{end}}=4+2+3+1=10\,\mathrm{ms}
\end{equation}
로 계산한다. 계약 레코드는 이 네 매개변수의 적용 환경과 증거 식별자를 가진다.
창 길이, 신호 군지연과 실행 반응시간은 각자의 필드로 유지한다.

\subsection{시계 보정 주기와 시간 오차}
시계 보정 직후의 오차 반경을 $\epsilon_0=0.2\,\mathrm{ms}$,
보정 간격 동안의 상대 속도 오차 상한을 $\rho$로 둔다.
잔차를 이 오차 반경에 포함하는 분석 모형은
\begin{equation}
 \epsilon(t)=\epsilon_0+\rho t
 \label{eq:drift_example}
\end{equation}
이다. $\rho=50\,\mathrm{ppm}$이면 $t=10\,\mathrm{s}$에서 $\epsilon=0.7\,\mathrm{ms}$,
$\rho=100\,\mathrm{ppm}$이면 같은 시각에서 $\epsilon=1.2\,\mathrm{ms}$를 얻는다.
양의 속도 오차 상한 $\rho>0$과 허용 반경 $\epsilon_\star$를 선택하면 보정 간격은
\begin{equation}
 T_{\mathrm{sync}}\le\frac{\epsilon_\star-\epsilon_0}{\rho},
 \qquad \epsilon_\star\ge\epsilon_0
\end{equation}
를 따른다. 그림~\ref{fig:clock}은 식~\eqref{eq:drift_example}의 계산값을 나타낸다.

\begin{figure}[tb]
\centering
\begin{tikzpicture}
\begin{axis}[width=.83\linewidth,height=58mm,
 xlabel={보정 이후 경과시간 (s)},ylabel={오차 반경 (ms)},
 xmin=0,xmax=20,ymin=0,ymax=2.4,xtick={0,5,10,15,20},ytick={0,.5,1,1.5,2},
 axis lines=left,tick align=outside,
 label style={font=\small},tick label style={font=\small},
 legend style={font=\small,draw=none,at={(.03,.97)},anchor=north west},
 legend cell align=left]
\addplot[black,solid,line width=.65pt,domain=0:20,samples=2]{0.2+0.05*x};
\addlegendentry{$\rho=50\,\mathrm{ppm}$}
\addplot[black,dashed,line width=.65pt,domain=0:20,samples=2]{0.2+0.1*x};
\addlegendentry{$\rho=100\,\mathrm{ppm}$}
\end{axis}
\end{tikzpicture}
\caption{시계 보정 간격과 오차 반경의 해석적 관계. 초기 반경 0.2 ms와 두 속도 오차 상한을 식~\eqref{eq:drift_example}에 대입하였다.}
\label{fig:clock}
\end{figure}

\subsection{고밀도 전기생리의 payload와 버퍼}
Neuropixels 1.0의 AP 스트림은 384채널을 $30\,\mathrm{kHz}$로 제공하고,
SpikeGLX의 아날로그 저장값은 signed 16-bit word를 사용한다\cite{neuropixels,spikeglx_metadata}.
이 AP 표본 배열의 payload 유량은
\begin{equation}
 r_{16}=384\times30\,000\times2=23.04\,\mathrm{MB/s}
\end{equation}
이다. 동일 표본 배열의 \code{f64} 표현은
\begin{equation}
 r_{64}=384\times30\,000\times8=92.16\,\mathrm{MB/s}
\end{equation}
를 사용한다. 여기서 MB는 $10^6$ byte이다. 각 값은 AP 표본 배열의 명시한 표현을 계량한다.
LFP, SY, 헤더와 재전송은 해당 경로의 추가 항으로 계량한다.

계산 모형은 기준 시계에서 균일한 $30\,000\,\mathrm{s}^{-1}$의 표본 발생을 사용한다.
한 블록을 채널당 30표본으로 구성하면 블록률은 $1000\,\mathrm{s}^{-1}$이고,
16-bit 블록의 크기는 $b=23\,040$ byte이다. 해당 FIFO 경계에
$R=50\,\mathrm{MB/s}$, $T=2\,\mathrm{ms}$의 서비스 곡선을 부여하면
\begin{equation}
 \begin{aligned}
 B_{\max}&=b+r_{16}T=69\,120\,\mathrm{byte}=67.5\,\mathrm{KiB},\\
 W_{\max}&=T+\frac bR=2.4608\,\mathrm{ms}.
 \end{aligned}
\end{equation}
블록 첫 표본부터 완성까지의 시간은 $29/30\,000\,\mathrm{s}$이다.
여기에 FIFO 지연을 합하면 첫 표본에서 선택한 출력 경계까지 $3.42747\,\mathrm{ms}$의
계산 상한을 얻는다. 서비스 계약은 블록 완성 경계를 기준으로 $R,T$를 지정한다.

연속 유입을 $T_{\mathrm{pause}}$ 동안 보관하는 값 배열의 용량은
$M=rT_{\mathrm{pause}}$로 계산한다. 그림~\ref{fig:memory}는 두 표현의 관계를 나타낸다.
\begin{figure}[tb]
\centering
\begin{tikzpicture}
\begin{axis}[width=.83\linewidth,height=58mm,
 xlabel={보관 시간 (ms)},ylabel={표본 배열 용량 (MB)},
 xmin=0,xmax=200,ymin=0,ymax=20,xtick={0,50,100,150,200},ytick={0,5,10,15,20},
 axis lines=left,tick align=outside,
 label style={font=\small},tick label style={font=\small},
 legend style={font=\small,draw=none,at={(.03,.97)},anchor=north west},
 legend cell align=left]
\addplot[black,solid,line width=.65pt,domain=0:200,samples=2]{23.04*x/1000};
\addlegendentry{16-bit: $23.04\,\mathrm{MB/s}$}
\addplot[black,dashed,line width=.65pt,domain=0:200,samples=2]{92.16*x/1000};
\addlegendentry{64-bit: $92.16\,\mathrm{MB/s}$}
\end{axis}
\end{tikzpicture}
\caption{384채널, 30 kHz AP 표본 배열의 보관 용량. 식 $M=rT_{\mathrm{pause}}$에 16-bit와 64-bit 표현의 유량을 대입하였다.}
\label{fig:memory}
\end{figure}

\subsection{EEG와 fNIRS의 시간 결합}
EEG 스트림은 전극 전압 스키마를, fNIRS 스트림은
$(s,d,\lambda,q,u)$의 광학 스키마를 유지한다. Aurora의 원시 광강도와 혈색소 변화량
출력은 서로 다른 명목 측정량을 갖는다\cite{nirx_formats}.
결합 연산자는 공통 시간 창 $[u_k,v_k)$와 각 입력의 시계 변환을 지정한다.
각 스트림에서 해당 창의 특징을 계산하고, 공통 frontier가 $v_k$에 도달하면
두 특징을 $(\varphi_{\mathrm{EEG}},\varphi_{\mathrm{fNIRS}})$의 곱 형식으로 출력한다.
후속 모델은 특징의 시간 지지, 표본 수, 유효성 마스크와 변환 계보를 함께 받는다.

\section{컴파일 산출물과 적합성 조건}
\label{sec:conformance}
\subsection{실행 manifest}
컴파일러는 프로그램의 정규화된 구문, 형식화 그래프, 하드웨어 바인딩과 실행 계획을
하나의 manifest로 생성한다. manifest는 값의 표현과 의미를 함께 추적하는 기준 문서이다.
표~\ref{tab:manifest}는 그 필드 집합을 정의한다.

\begin{table}[tb]
\centering\small
\caption{실행 manifest의 필수 구성}
\begin{tabularx}{\linewidth}{@{}P{.23\linewidth}Y@{}}
\toprule 구성 & 저장 내용 \\\midrule
언어와 그래프 & 언어 버전, 정규화 AST, 그래프 식별자, 원시 연산자 버전 \\
하드웨어 & 프로필 $H$, 설정 $\theta$, 기능 요구, 접근 조건, 원문 식별자 \\
형식 & 채널 순서, 측정량, 단위, reference, 좌표계, 자료형, stride \\
시간 & 원래 시계, clock epoch, 변환 버전, 유효 구간, 오차 계약 \\
수치 & 계수, 표현 정밀도, 반올림ㆍ포화 정책, 참조 오차 기준 \\
학습 & 분할 manifest, fit 계보, 모델 digest, 모델 활성화 사건 \\
실행 & actor 소유권, 일정, 버퍼 용량, 큐 정책, 서비스ㆍ응답시간 계약 \\
기록 & 원본 범위, 품질 사건, 스키마 변경, 실행 환경과 결과 식별자 \\
\bottomrule
\end{tabularx}
\label{tab:manifest}
\end{table}

형식화 구체화는 단위ㆍ채널ㆍ시계 제약을 풀고 각 연결선의 타입을 결정한다.
정적 인덱스 제약은 정수 선형 산술, 유리수 비교와 유한 집합 관계로 구성한다.
행렬의 수치 조건과 장치 조회값은 해당 원시 연산자ㆍ프로필의 계약 식별자에 연결한다.
원시 연산자 선택은 수치 커널과 상태 표현을 결정한다. 실행 계획 생성은
고정률 부분그래프의 firing 순서, 사건 경계, 메모리 소유권과 버퍼 위치를 결정한다.
각 변환은 입력ㆍ출력 그래프의 대응표와 계약 증거를 manifest에 추가한다.

\subsection{관측 의미의 대응}
프로그램 $P$, 입력 trace $I$, 환경 $A$에서 코어 의미가 허용하는 실행 집합을
$\mathcal T(P,I,A)$로 둔다. 정수 디코딩, 사건 순서와 모델 버전은 정확한 관측 관계를
사용하고, 실수 결과는 식~\eqref{eq:numeric}의 수치 관계를 사용한다.
컴파일 결과 $C(P)$의 실행 trace $T$에 대해 적합성 관계를
\begin{equation}
 \operatorname{Obs}(T)
 \in\mathcal T_{\epsilon}(P,I,A)
 \label{eq:refinement}
\end{equation}
로 정의한다. $\mathcal T_\epsilon$은 명시한 수치 허용오차와 시간 구간을 적용한 관측 집합이다.
관측에는 값, 가용 시각, 품질, 모델 버전, schemaㆍclock epoch와 외부 사건을 포함한다.
버퍼 정책이 생성하는 gap과 상태 전이도 같은 trace의 원소로 기록한다.

검사기는 각 원시 연산자의 전이 불변식을 확인하고, 합성 규칙에 따라 출력 계약을 계산한다.
하드웨어 바인딩은 식~\eqref{eq:bind}를 평가한다. 일정 검사는 반복 벡터와 실행 접두의
토큰 수를 검사하고, 시간 검사는 선택한 arrivalㆍservice curve 또는 응답시간 계약을 계산한다.
각 결과는 적용한 입력 영역, 구성, 도구ㆍ계약 버전과 함께 저장한다.

\subsection{규격 적합성의 판정 대상}
\begin{longtable}{@{}P{.19\linewidth}P{.36\linewidth}P{.39\linewidth}@{}}
\caption{적합성 검사 대상과 판정 관계}\label{tab:conformance}\\
\toprule 대상 & 입력과 근거 & 판정 관계 \\\midrule
\endfirsthead
\multicolumn{3}{l}{\small 표 \thetable. 적합성 검사 대상과 판정 관계 (계속)}\\
\toprule 대상 & 입력과 근거 & 판정 관계 \\\midrule
\endhead
\bottomrule\endfoot
패킷 디코딩 & 프로토콜 판본, 패킷ㆍ기준값 쌍, 경계 길이 & 부호ㆍbyte orderㆍ채널ㆍ상태값의 정확한 대응 \\
물리량 변환 & 정수 범위, gain, offset, 기준 단위 & 기준값과 식으로 정의한 수치 관계 \\
시계 변환 & 동기 사건 쌍, 관측 경계, 유효 구간 & 시간 구간 포함과 오차 증거 종류 \\
스키마 전이 & 설정 전후 descriptor와 경계 사건 & 구간별 형식, 상태 변환, epoch 순서 \\
품질 전파 & gapㆍ보간ㆍ복구ㆍ포화 사건 trace & 연산자 정책에 따른 출력 마스크와 계보 \\
학습 계보 & 원천 ID, 그룹 함수, fitㆍtuneㆍeval 집합 & 식~\eqref{eq:split} 및 모델 가용 조건 \\
실행 일정 & 생산ㆍ소비율, 반복 벡터, 초기 토큰 & firing 가능성, 주기 평형, 최대 점유량 \\
시간ㆍ메모리 & 유입 상한, 서비스 곡선, 환경 조건 & backlogㆍ반응시간ㆍ메모리 예산 \\
기록ㆍ재생 & manifest, 원본ㆍ사건 순서, 활성 버전 & 식~\eqref{eq:refinement}의 관측 대응 \\
\end{longtable}

\subsection{설정 변경과 재연결의 전이}
설정 변경 사건은 새 구성 $\theta'$와 새 스키마 기술자를 생성한다.
실행기는 현재 epoch의 진행 경계를 닫고, 남은 처리 의존성을 해당 epoch에 귀속시킨다.
다음 epoch에는 새로운 형식과 호환하는 상태 변환 $R_Z$를 적용한다.
\begin{equation}
 R_Z:Z_{\theta}\longrightarrow Z_{\theta'},\qquad
 R_Z(\mathcal J_{\theta})\subseteq\mathcal J_{\theta'}.
\end{equation}
초기화 정책은 $R_Z(z)=z^0_{\theta'}$로 표현한다.
장치 재연결은 acquisition epoch를 갱신하고, 시계 재설정은 clock epoch와
새 시간 변환 레코드를 생성한다. 이 사건들이 배열 해석과 시간 연속성의 경계를 결정한다.

\subsection{기록과 재생의 관측 경계}
기록 어댑터는 원래 측정 시각, 원래 가용 시각과 재생 시각을 각각 저장한다.
실시간 실행 재생은 원래 가용 순서를 사용하여 $\cut_u$를 구성한다.
회고 분석은 새 시간 뷰와 변환 계보를 생성하여 기록 구간을 입력으로 사용한다.
XDF는 스트림 timestamp와 clock offset 정보를 기록하고\cite{xdf,lsl_time},
SNIRF는 광학 측정열ㆍ시간 배열과 chunked 기록을 표현한다\cite{snirf}.
Open Ephys와 SpikeGLX의 기록 어댑터는 해당 버전의 표본 번호ㆍ시각ㆍ스케일 필드를
복원한다\cite{open_ephys_binary,spikeglx_metadata}.
재생기는 이 기록 형식들을 공통 시간 레코드와 계보 그래프에 대응시킨다.

\section{설계의 종합}
본 설계는 BCI 프로그램을 물리량ㆍ채널ㆍ시간ㆍ계보가 붙은 스트림의 합성으로 정의한다.
원시 연산자의 상태 불변식과 가용 시각 인과성은 직렬ㆍ병렬 합성의 의미를 구성하고,
초기 지연은 폐루프 그래프의 유일한 실행을 결정한다.
학습 효과와 모델 활성화 사건은 데이터 분할, 전처리 통계, 예측 시점의 매개변수를 연결한다.

하드웨어 계층은 지원 프로필, 유효 설정 공간과 구성별 기능 집합을 사용한다.
시계 변환은 유효 구간과 오차 표현을 제공하며, 창 연산은 pending 상태를 포함한
진행 경계로 완결 시점을 결정한다. Rust 실행 계획은 가변 상태의 소유권,
불변 블록의 공유, 고정률 일정과 서비스 계약을 결합한다.
컴파일 manifest와 실행 trace는 설계의 형식 조건을 실제 장치 구성과 연산 결과에 대응시킨다.
이 구조에서 장치 확장은 프로필과 어댑터 계약을, 알고리즘 확장은 원시 연산자와 학습 효과를 추가한다.

\FloatBarrier
\appendix
\section{하드웨어 지원 프로필}\label{app:hardware}

표~\ref{tab:hardware_profiles}는 장치와 획득 경로를 결합한 18개 지원 프로필을 정리한다.
각 프로필은 모델, 펌웨어, 획득 모드와 소프트웨어 판본을 식별하고,
채널 구성과 표본 표현을 해당 설정에 연결한다.
표의 필드는 본 논문의 프로필 스키마를 따르며,
제조사 규격, 장치 조회 결과와 획득 프로그램의 설정을 값의 출처로 기록한다.

\begingroup
\small
\begin{longtable}{@{}P{.16\linewidth}P{.35\linewidth}P{.43\linewidth}@{}}
\caption{장치 및 획득 경로별 지원 프로필}\label{tab:hardware_profiles}\\
\toprule
식별자와 대상 & 획득 경로 및 신호 규격 & 주요 필드와 값의 출처 \\
\midrule
\endfirsthead
\caption[]{장치 및 획득 경로별 지원 프로필 (계속)}\\
\toprule
식별자와 대상 & 획득 경로 및 신호 규격 & 주요 필드와 값의 출처 \\
\midrule
\endhead
\midrule
\multicolumn{3}{r}{다음 쪽에 계속}\\
\endfoot
\bottomrule
\endlastfoot

A01\newline OpenBCI\newline Cyton
& 전용 무선 동글과 USB serial로 연결한다. 기본 모드는 8 EEG 채널을 250 SPS로 전송한다. 33바이트 프레임에 채널별 signed 24-bit big-endian 값과 보조 데이터를 배치한다.\citep{cyton}
& \code{frame}, \code{counter}, \code{aux\_mode}, \code{gain}을 기록한다. 공식 프레임 규격과 보드 설정 응답에서 배열, 계수기 폭, 보조부 종류와 전압 환산 배율을 얻는다. \\
\addlinespace[5pt]

A02\newline Cyton+Daisy
& 두 8채널 bank가 무선 경로를 공유한다. ADC는 채널당 250 SPS로 획득하고, 각 bank는 연속 두 표본의 평균을 125 SPS로 교대 전송한다.\citep{cyton}
& \code{adc\_rate}, \code{output\_rate}, \code{bank\_phase}, \code{aggregation}을 기록한다. 공식 Daisy 전송 규칙으로 평균 구간과 교대 위상을 정하고 첫 패킷 뒤에서 유효 구간을 시작한다. \\
\addlinespace[5pt]

A03\newline OpenBCI\newline Ganglion
& BLE로 4 EEG 채널을 200 Hz로 전달한다. 압축 신호 패킷은 20바이트에 두 표본을 담는다. 펌웨어 2 계열은 delta 압축을, 3.0.0부터는 18/19-bit MSB 전송을 사용한다.\citep{ganglion,ganglion_release}
& \code{firmware}, \code{codec}, \code{packet\_id\_domain}, \code{connection\_epoch}를 기록한다. 공식 펌웨어 tag에서 복원 규칙과 패킷 식별자 범위를 정하고 접속 재설정을 획득 경계에 반영한다. \\
\addlinespace[5pt]

A04\newline actiCHamp\newline 및 Plus
& Amplifier SDK로 획득한다. 32 EEG+8 AUX 구성은 최대 100 kHz, 64+8은 50 kHz, 96/128/160+8은 25 kHz를 제공한다. 채널별 ADC는 24-bit로 변환한다.\citep{actichamp,brainproducts_sdk}
& \code{channel\_modules}, \code{sampling\_mode}, \code{aux\_count}, \code{sdk\_version}을 기록한다. 공식 채널--속도 표가 허용 조합을 정하고 SDK 조회 결과가 활성 구성을 채운다. \\
\addlinespace[5pt]

A05\newline Brain Products\newline RDA
& BrainVision Recorder가 설정한 EEG 채널과 사건 정보를 RDA/TCP 경로로 전달한다. 장치 선택, 획득 표본률과 출력 구성은 해당 Recorder 세션을 따른다.\citep{recorder}
& \code{recorder\_version}, \code{amplifier\_model}, \code{output\_schema}, \code{marker\_layout}을 기록한다. Recorder 작업 설정과 해당 판본의 스트림 기술에서 채널 배열 및 사건 표현을 얻는다. \\
\addlinespace[5pt]

A06\newline Brain Products\newline LSL
& 제조사 connector가 actiCHamp, BrainAmp 또는 LiveAmp 획득값을 LSL로 제공한다. 장치와 connector 설정이 채널 구성, 표본률과 기준 전위의 의미를 결정한다.\citep{brainproducts_sdk}
& \code{connector\_version}, \code{device\_driver}, \code{reference}, \code{timestamp\_boundary}를 기록한다. connector 문서와 LSL 기술자를 연결하여 신호 유도 방식과 시각 관측 위치를 명시한다. \\
\addlinespace[5pt]

A07\newline Unicorn\newline Hybrid Black
& Bluetooth 2.1+EDR와 C API를 사용한다. 8 EEG 채널을 250 Hz로 획득하며, API는 활성 채널의 float32 값을 표본 순서로 배열한다. 관성 및 상태 채널도 함께 구성한다.\citep{unicorn_manual,unicorn_api}
& \code{active\_channels}, \code{scan\_stride}, \code{counter\_channel}, \code{validation\_channel}을 기록한다. API 채널 조회가 신호, 계수기와 유효성 지시자의 배열 위치를 결정한다. \\
\addlinespace[5pt]

A08\newline g.USBamp\newline RESEARCH
& USB 2.0과 g.NEEDaccess 경로를 사용한다. 16 monopolar 또는 8 bipolar 채널에서 24-bit 동시 표본화를 제공하며, 채널당 표본률 상한은 38.4 kHz이다.\citep{gusbamp}
& \code{derivation\_mode}, \code{sampling\_mode}, \code{dsp\_filters}, \code{server\_version}을 기록한다. 제품 규격과 서버 설정 조회에서 유도 방식, 출력률과 적용 필터를 얻는다. \\
\addlinespace[5pt]

A09\newline ANT eego\newline EE-2XX/4XX
& eego SDK와 BrainFlow 경로를 연결한다. 제조사 연동 문서는 EE-2XX/EE-4XX와 8채널 EE-411 예제를 다룬다. 선택 모델의 SDK 설정이 채널 수와 획득률을 정한다.\citep{ant_eego}
& \code{device\_model}, \code{sdk\_package}, \code{channel\_map}, \code{timestamp\_origin}을 기록한다. SDK 조회와 bridge 기술자에서 채널 단위, 배열과 시각의 생성 위치를 얻는다. \\
\addlinespace[5pt]

A10\newline BioSemi\newline ActiveTwo/2.5
& USB 수신기의 ringbuffer SDK로 획득한다. 일반 모드 4/5/6/7은 각각 2048/4096/8192/16384 Hz를 사용한다. ADC값은 32-bit word의 상위 24비트에 놓인다.\citep{biosemi_sdk,biosemi_rate}
& \code{speedmode}, \code{words\_per\_sample}, \code{status\_layout}, \code{integer\_scale}을 기록한다. Mk2의 해당 모드별 프레임은 282/154/90/58 word이며 공식 buffer 표가 상태부와 환산식을 정한다. \\
\addlinespace[5pt]

A11\newline BioSemi\newline ActiveThree
& 장치는 설치 채널을 16384 Hz로 전송하고 소프트웨어가 저장 및 분석 표본률을 조정한다. 수신 SDK는 ActiveThree 전용 ringbuffer 배열을 사용한다.\citep{biosemi_rate,biosemi_a3}
& \code{schema\_id}, \code{wire\_rate}, \code{output\_rate}, \code{decimation}을 기록한다. 해당 SDK 판본에서 프레임 배열을 얻고 획득 설정에서 출력률과 표본 축소 필터를 얻는다. \\
\addlinespace[5pt]

A12\newline Intan RHD\newline Rhythm USB-7310
& SPI headstage와 USB 3.0 경로를 사용한다. XEM7310 구성은 최대 16개 RHD chip, 1024채널을 제공하며 채널당 1--30 kS/s 범위의 지원률을 선택한다.\citep{intan_rhd}
& \code{chip\_map}, \code{sampling\_rate}, \code{fpga\_image}, \code{api\_version}을 기록한다. Rhythm 3.3 API와 FPGA 설정에서 활성 스트림, 표본률 열거값과 프레임 배치를 얻는다. \\
\addlinespace[5pt]

A13\newline Intan RHS\newline RhythmStim
& RhythmStim USB-7310은 최대 8개 RHS2116 chip, 128채널을 연결한다. 채널당 20/25/30 kS/s 중 하나를 선택하며 기록 신호와 자극 상태를 전달한다.\citep{intan_rhs}
& \code{chip\_map}, \code{sampling\_rate}, \code{stim\_state}, \code{trigger\_mode}를 기록한다. RhythmStim 3.3 문서와 장치 설정에서 채널 배치, 자극 상태와 사건 구동 방식을 얻는다. \\
\addlinespace[5pt]

A14\newline Intan\newline RHX/TCP
& RHX가 RHD/RHS의 선택 파형을 TCP로 전달하고 명령 인터페이스로 획득을 제어한다. 활성 장치, 채널 선택과 실시간 처리 설정이 출력 신호를 구성한다.\citep{intan_rhx}
& \code{rhx\_version}, \code{channel\_selection}, \code{filter\_chain}, \code{reference}를 기록한다. RHX 설정과 TCP 인터페이스 문서에 출력 배열, 필터 순서와 기준 전위를 연결한다. \\
\addlinespace[5pt]

A15\newline Neuropixels 1.0\newline SpikeGLX
& probe당 AP 384채널을 30 kHz로, LFP 384채널을 2.5 kHz로 획득한다. SpikeGLX가 probe와 보조 신호를 통합하고 원격 접근 경로를 제공한다.\citep{neuropixels,spikeglx}
& \code{probe\_type}, \code{electrode\_map}, \code{band\_rates}, \code{gain}을 기록한다. probe 메타데이터와 획득 설정에서 전극 선택, 대역별 시간 격자와 신호 환산값을 얻는다. \\
\addlinespace[5pt]

A16\newline Blackrock\newline CereLink
& CereLink는 장치 UDP, 다른 CereLink의 공유메모리, Central 공유메모리를 연결 모드로 제공한다. 장치 모델과 sample group 설정이 채널 배열 및 출력률을 결정한다.\citep{cerelink}
& \code{connection\_mode}, \code{hardware\_protocol}, \code{sample\_group}, \code{shmem\_schema}를 기록한다. 프로젝트의 protocol 정의와 장치 응답에서 전송 및 공유메모리 자료형을 얻는다. \\
\addlinespace[5pt]

A17\newline NIRSport2\newline Aurora/LSL
& 단일 장치는 광원 8/16개, 검출기 8--16개와 760/850 nm 파장을 사용한다. 단일ㆍ다장비ㆍBorealis 모드의 제품 표는 각각 최대 240 Hz를 제시한다.\citep{nirsport,nirx_formats}
& \code{channel\_tuple}, \code{actual\_rate}, \code{quantity\_kind}, \code{montage}를 기록한다. Aurora header와 LSL 기술자에서 광원--검출기--파장 조합, 실제 출력률과 광학 측정량을 얻는다. \\
\addlinespace[5pt]

A18\newline Artinis\newline OxySoft/LSL
& OxySoft가 선택 장치의 데이터를 획득하고 혈색소 농도 변화와 조직 포화도를 계산하며 LSL 경로를 제공한다. 장치 구성과 내보내기 설정이 채널 및 표본률을 정한다.\citep{oxysoft}
& \code{device\_model}, \code{oxysoft\_version}, \code{quantity\_kind}, \code{processing\_parameters}를 기록한다. 설정과 스트림 기술자에서 측정량, 단위와 적용 처리값을 얻는다. \\

\end{longtable}
\endgroup

\clearpage
\section{기호와 manifest 식별자}
\label{app:symbols}
\begin{longtable}{@{}P{.20\linewidth}P{.74\linewidth}@{}}
\caption{주요 기호의 의미}\\
\toprule 기호 & 의미 \\\midrule
\endfirsthead
\toprule 기호 & 의미 \\\midrule
\endhead
\bottomrule\endfoot
$D$ & 표본의 수치 표현 \\
$T_{\mathrm{val}}$ & 벡터ㆍepochㆍ행렬ㆍ사건을 포함한 값 형식 \\
$\mathcal C,\chi_c$ & 순서가 있는 채널 스키마와 채널 기술자 \\
$\kappa,\mathcal R$ & 시계 영역과 표본 발생 규칙 \\
$X\in D^{C\times N}$ & 채널 $C$개, 표본 $N$개의 값 블록 \\
$I_i^{\mathrm{meas}}$ & 대표 측정 시각의 구간 \\
$a_i$ & 지정한 관측 경계에서의 가용 시각 \\
$\cut_u(S)$ & 시각 $u$까지 가용한 스트림 부분 \\
$\rho_i,\rho_v^{\mathrm{fit}}$ & 레코드 계보와 모델 학습 계보 \\
$e^{\mathrm{schema}},e^{\mathrm{clock}}$ & 스키마와 시계의 세대 식별자 \\
$H,\theta,\Omega_H$ & 하드웨어 프로필, 선택 구성, 유효 구성 공간 \\
$C(H,\theta),R_P$ & 구성별 기능 집합과 프로그램의 기능 요구 \\
$Z_o,\mathcal J_o,\delta_o$ & 연산자 상태 공간, 상태 불변식, 상태 전이 \\
$\mathcal E_o$ & 처리 모드, 요구 기능 집합, 전방 관측 길이 \\
$M_v,\vartheta_v$ & 모델 버전과 수치 매개변수 \\
$\alpha,\beta$ & 누적 유입과 서비스의 경계 함수 \\
$b,r,R,T$ & burst, 지속 유입률, 서비스율, 서비스 지연 \\
\end{longtable}

manifest의 모든 식별자는 이름 공간과 버전을 함께 가진다.
장치ㆍ수신기 식별자, 원천 스트림 식별자, 획득 epoch, 데이터 블록 순번은
원천 레코드의 범위를 결정한다. 스키마ㆍ시간 뷰ㆍ모델ㆍ계약 식별자는
변환 관계를 결정한다. 수치 상수와 설정 레코드는 정규화한 표현의 digest를 사용한다.
기록기는 이 식별자와 전이 사건을 함께 저장하여 동일 원천에서 생성한 여러 분석 결과를 연결한다.

\FloatBarrier
\begingroup
\small
\raggedright
\interlinepenalty=10000
\setstretch{1.08}
\setlength{\bibsep}{4pt}
\bibliographystyle{unsrtnat}
\bibliography{references}
\endgroup
\end{document}
